% =====================================================================
%  Rapport de stage — Campus Navigator (ENSAM Rabat)
%  Compilation : pdflatex rapport.tex (deux passes pour la table des matières)
%  Les captures d'écran sont attendues dans le dossier images/ ;
%  une image absente est remplacée par un cadre portant son nom.
% =====================================================================
\documentclass[12pt,a4paper]{report}

\usepackage[utf8]{inputenc}
\usepackage[T1]{fontenc}
\usepackage{lmodern}
\usepackage[french]{babel}
\usepackage[margin=2.5cm]{geometry}
\usepackage{graphicx}
\usepackage{amsmath,amssymb}
\usepackage{array,booktabs,tabularx,longtable}
\usepackage{xcolor}
\usepackage{tikz}
\usetikzlibrary{arrows.meta,positioning,calc}
\usepackage{pgfgantt}
\usepackage{listings}
\usepackage{caption}
\usepackage{subcaption}
\usepackage{titlesec}
\usepackage{enumitem}
\usepackage{chngcntr}
\usepackage{url}
\usepackage[hidelinks]{hyperref}

\setlength{\parindent}{1.5em}
\setlength{\parskip}{0.3em}
\renewcommand{\baselinestretch}{1.15}

% Numérotation des figures et tableaux par section (ex. Figure 1.7.1)
\counterwithin{figure}{section}
\counterwithin{table}{section}

\titleformat{\chapter}[hang]{\Huge\bfseries}{\thechapter\hspace{0.8em}}{0pt}{\Huge\bfseries}
\titlespacing*{\chapter}{0pt}{60pt}{40pt}

\definecolor{codebg}{RGB}{246,248,250}
\lstset{
  basicstyle=\ttfamily\small,
  backgroundcolor=\color{codebg},
  frame=single, rulecolor=\color{black!20},
  breaklines=true, columns=fullflexible,
  showstringspaces=false,
  literate={é}{{\'e}}1 {è}{{\`e}}1 {à}{{\`a}}1 {ê}{{\^e}}1 {É}{{\'E}}1 {ô}{{\^o}}1 {ç}{{\c c}}1
}

% Image avec repli : si le fichier n'existe pas, un cadre indique le nom attendu
\newcommand{\img}[2]{%
  \IfFileExists{images/#2}%
    {\includegraphics[width=#1]{images/#2}}%
    {\fbox{\parbox[c][4cm][c]{0.9\linewidth}{\centering\small\ttfamily Image à insérer :\\ images/\detokenize{#2}}}}%
}
\emergencystretch=3em

\begin{document}

% =====================================================================
% Page de garde
% =====================================================================
\begin{titlepage}
  \centering
  \IfFileExists{public/logo-ensam.jpg}{\includegraphics[width=0.45\linewidth]{public/logo-ensam.jpg}\par}{}
  \vspace*{2cm}
  {\Huge\bfseries Rapport de Stage\par}
  \vspace{1.5cm}
  {\large\bfseries Effectué au sein de :\par}
  \vspace{0.2cm}
  {\Large ENSAM de Rabat\par}
  \vspace{2cm}
  {\large\bfseries Effectué :\par}
  {\large Du 8 Juillet 2026 au 30 Août 2026\par}
  \vspace{2.5cm}
  {\large\bfseries Réalisé par :\par}
  \vspace{0.4cm}
  {\large BAHLOULI Zaineb \hfill OUAZZANI TAYBI Walae\par}
  \vspace{1.5cm}
  {\large\bfseries Encadré par :\par}
  \vspace{0.2cm}
  {\large PR. Alami Kamouri Sophia\par}
  \vfill
  {\large Octobre 2026\par}
\end{titlepage}

\pagenumbering{Roman}

% =====================================================================
% Résumé
% =====================================================================
\section*{Résumé}
\addcontentsline{toc}{chapter}{Résumé}

Ce projet porte sur la conception et le développement de \emph{Campus Navigator}, un système
de navigation intelligent destiné au campus de l'ENSAM Rabat. Il a pour objectif d'aider les
étudiants, les enseignants et les visiteurs à trouver facilement leur chemin vers les différents
lieux du campus (salles de cours, laboratoires, bureaux administratifs, amphithéâtres, salle
d'étude\ldots) grâce à une interface conversationnelle bilingue (français et anglais).

Le système repose sur une architecture client-serveur. Un backend FastAPI héberge le graphe du
campus (modélisé à partir d'un classeur Excel), identifie la destination demandée à l'aide d'un
modèle de langage interrogé via l'API Groq, complété par un mécanisme de recherche locale
approximative lorsque le modèle n'est pas disponible, calcule le plus court chemin avec
l'algorithme de Dijkstra (NetworkX) et génère une instruction de marche par point de passage.
Un frontend React assure le guidage en temps réel : le trajet restant est tracé sur les plans SVG
réels du campus, la carte change automatiquement d'étage et suit l'utilisateur, dont la position
est estimée par le GPS du téléphone à l'extérieur et par un podomètre exploitant l'accéléromètre à
l'intérieur. En cas d'écart au trajet, l'itinéraire est recalculé automatiquement. L'application
est conteneurisée avec Docker et déployée sur une adresse HTTPS publique, accessible par simple
scan d'un QR Code depuis n'importe quel réseau.

\medskip
\noindent\textbf{Mots-clés :} navigation intérieure, théorie des graphes, plus court chemin,
traitement du langage naturel, géolocalisation, podométrie, expérience utilisateur.

\clearpage
% =====================================================================
% Abstract
% =====================================================================
\section*{Abstract}
\addcontentsline{toc}{chapter}{Abstract}

This project focuses on the design and development of \emph{Campus Navigator}, an intelligent
navigation system for the ENSAM Rabat campus. Its purpose is to help students, professors and
visitors easily find their way to campus locations (classrooms, laboratories, administrative
offices, lecture halls, study room\ldots) through a bilingual (French and English)
conversational interface.

The system follows a client-server architecture. A FastAPI backend hosts the campus graph
(built from an Excel workbook), identifies the requested destination with a large language
model served by the Groq API, backed by a local fuzzy-matching fallback when the model is
unavailable, computes the shortest path with Dijkstra's algorithm (NetworkX) and generates one
walking instruction per waypoint. A React frontend provides real-time guidance: the remaining
route is drawn on the real SVG floor plans, the map switches floors automatically and follows
the user, whose position is estimated with the phone's GPS outdoors and with an
accelerometer-based step counter indoors. When the user leaves the route, it is recomputed
automatically. The application is containerised with Docker and deployed on a public HTTPS
address, reachable by scanning a QR code from any network.

\medskip
\noindent\textbf{Keywords:} indoor navigation, graph theory, shortest path, natural language
processing, geolocation, pedestrian dead reckoning, user experience.

% =====================================================================
% Abréviations
% =====================================================================
\chapter*{Liste des abréviations}
\addcontentsline{toc}{chapter}{Liste des abréviations}

\begin{description}[leftmargin=1.5em, itemsep=0.25em]
  \item[API] : \emph{Application Programming Interface}, interface de programmation applicative
  \item[BLE] : \emph{Bluetooth Low Energy}, Bluetooth à basse consommation
  \item[CORS] : \emph{Cross-Origin Resource Sharing}, partage de ressources entre origines
  \item[DAX] : \emph{Data Analysis Expressions}
  \item[ENSAM] : École Nationale Supérieure d'Arts et Métiers
  \item[GNSS] : \emph{Global Navigation Satellite System}, système mondial de navigation par satellites
  \item[GPS] : \emph{Global Positioning System}, système de positionnement global
  \item[HTTP(S)] : \emph{HyperText Transfer Protocol (Secure)}
  \item[IA] : Intelligence artificielle
  \item[INDIA] : Ingénierie Numérique, Data Science et Intelligence Artificielle
  \item[JSON] : \emph{JavaScript Object Notation}
  \item[LLM] : \emph{Large Language Model}, grand modèle de langage
  \item[PDR] : \emph{Pedestrian Dead Reckoning}, navigation pédestre à l'estime
  \item[QR Code] : \emph{Quick Response Code}
  \item[RDC] : Rez-de-chaussée
  \item[REST] : \emph{Representational State Transfer}
  \item[SIG] : Système d'information géographique
  \item[SVG] : \emph{Scalable Vector Graphics}
  \item[TALN] : Traitement automatique du langage naturel
  \item[URL] : \emph{Uniform Resource Locator}
  \item[UWB] : \emph{Ultra-Wideband}, ultra-large bande
  \item[UX] : \emph{User Experience}, expérience utilisateur
\end{description}

\listoffigures
\addcontentsline{toc}{chapter}{Table des figures}
\listoftables
\addcontentsline{toc}{chapter}{Liste des tableaux}
\tableofcontents

\clearpage
\pagenumbering{arabic}

% =====================================================================
% Introduction générale
% =====================================================================
\chapter*{Introduction Générale}
\addcontentsline{toc}{chapter}{Introduction Générale}

Les campus universitaires modernes sont souvent constitués de plusieurs bâtiments, étages et
espaces extérieurs dont l'agencement n'est pas toujours intuitif pour les nouveaux arrivants —
étudiants en première année, enseignants vacataires ou visiteurs externes. À l'ENSAM Rabat, cette
complexité se traduit concrètement par des salles réparties sur plusieurs bâtiments et niveaux,
reliées par des couloirs, des escaliers et des allées extérieures dont la disposition n'est pas
évidente au premier abord. Si les solutions de navigation extérieure (GPS, cartes routières) sont
aujourd'hui matures, la navigation à l'intérieur des bâtiments reste un problème nettement moins
bien résolu : le signal GPS y est peu fiable, et les plans papier affichés à l'entrée des bâtiments
offrent une expérience statique, peu adaptée à la recherche précise d'un lieu.

Comment permettre à un utilisateur de trouver rapidement et intuitivement le chemin vers
n'importe quel lieu du campus sans dépendre d'une signalétique physique figée ni de compétences
techniques particulières pour interroger le système ? La problématique se décompose en plusieurs
sous-questions techniques : comment modéliser le campus (bâtiments, étages, couloirs, allées)
sous une forme exploitable algorithmiquement ? Comment interpréter une requête formulée en
langage courant (« Comment aller à la salle B7 ? ») pour en extraire la destination réelle, y
compris lorsque le service d'IA est indisponible ? Comment calculer et représenter un itinéraire
cohérent lorsqu'il traverse plusieurs bâtiments et plusieurs étages ? Comment suivre l'utilisateur
pendant qu'il marche, sans infrastructure de localisation dédiée ? Et comment rendre le service
accessible depuis n'importe quel smartphone, quel que soit le réseau utilisé ?

Pour répondre à cette problématique, le projet a été mené selon une démarche progressive : une
phase de modélisation du campus sous forme de graphe pondéré (nœuds et arêtes, rattachés à un
étage et à un bâtiment) ; une phase de développement backend, exposant ce graphe via une API REST
(FastAPI), intégrant un modèle de langage (Groq) pour l'extraction de la destination et la
rédaction des instructions, et calculant le plus court chemin par l'algorithme de Dijkstra
(NetworkX) ; une phase de développement frontend, avec une interface de chat en React affichant
le trajet sur les plans réels de chaque étage ; une phase de calibration, de débogage et
d'optimisation de l'expérience mobile ; et enfin une phase d'extension consacrée au guidage en
temps réel (GPS à l'extérieur, podomètre à l'intérieur, recalcul automatique de l'itinéraire) et
au déploiement public de l'application.

Ce rapport s'articule en quatre chapitres. Le premier chapitre présente le cadre général du
projet : le contexte académique, la problématique et la planification adoptée. Le deuxième
chapitre dresse un état de l'art des technologies et outils mobilisés. Le troisième chapitre
détaille la méthodologie suivie, depuis la mise en place de l'environnement de développement
jusqu'à la conception de la base de données, du pipeline applicatif, du module de navigation en
temps réel et du déploiement. Le quatrième chapitre présente les résultats obtenus ainsi que les
principales difficultés rencontrées et les solutions apportées. Une conclusion générale dresse
enfin le bilan du travail réalisé et propose des perspectives d'amélioration.

% =====================================================================
% Chapitre 1
% =====================================================================
\chapter{Cadre général du projet}

\section{Introduction}

Dans un établissement universitaire, l'orientation des étudiants, des enseignants, du personnel
administratif et des visiteurs constitue un besoin important, particulièrement lorsque les
bâtiments comportent plusieurs niveaux et un nombre important de salles, de bureaux, de
laboratoires et de services. Pour une personne qui ne connaît pas suffisamment l'organisation
spatiale de l'établissement, retrouver un emplacement précis peut nécessiter du temps et
plusieurs demandes d'information.

Cette difficulté est particulièrement présente dans le contexte de la navigation intérieure.
Contrairement à la navigation en extérieur, où les systèmes de positionnement par satellite
permettent généralement de déterminer la position d'un utilisateur, la navigation à l'intérieur
des bâtiments nécessite de prendre en considération la structure des espaces, les couloirs, les
différents niveaux ainsi que les connexions entre les zones accessibles.

Dans ce contexte, le projet \textbf{Campus Navigator} a pour objectif de proposer une solution
numérique destinée à faciliter l'orientation et la navigation au sein de l'\textbf{École
Nationale Supérieure d'Arts et Métiers de Rabat (ENSAM de Rabat)}. La solution permet à
l'utilisateur de rechercher une destination, d'obtenir un itinéraire depuis sa position de
départ, de le visualiser directement sur le plan du campus et d'être guidé pas à pas pendant son
déplacement.

Le projet combine plusieurs domaines de l'ingénierie informatique, notamment le développement
web, la modélisation des données, la théorie des graphes, la cartographie numérique, la
géolocalisation et l'intelligence artificielle. L'objectif de ce premier chapitre est de
présenter le cadre général dans lequel le projet a été réalisé, d'exposer la problématique ayant
motivé son développement, de présenter les objectifs et le fonctionnement général de la
solution, puis de décrire les différentes étapes de sa réalisation.

\section{Présentation du cadre d'accueil}

\subsection{Présentation de l'établissement}

Le projet Campus Navigator a été réalisé au sein de l'\textbf{École Nationale Supérieure d'Arts
et Métiers de Rabat (ENSAM de Rabat)}. L'établissement constitue un environnement universitaire
regroupant différents espaces dédiés à l'enseignement, à la recherche, à l'administration et aux
activités de la communauté universitaire.

La diversité des espaces présents dans l'établissement implique une organisation spatiale
composée de plusieurs bâtiments (bâtiment administratif, atelier, bâtiment de l'amphithéâtre,
département électrique, amphithéâtres, salle d'étude, services\ldots), de salles, de bureaux,
de laboratoires, de zones de circulation et d'allées extérieures. Cette organisation peut rendre
les déplacements difficiles pour les personnes qui connaissent peu le campus.

Dans ce contexte, la mise en place d'un système numérique permettant de rechercher un lieu et
d'obtenir un itinéraire représente une approche permettant de faciliter l'accès à l'information
spatiale. Campus Navigator s'inscrit ainsi dans une démarche visant à exploiter les technologies
numériques pour améliorer l'orientation au sein de l'établissement.

\subsection{Présentation de la filière}

Le projet a été réalisé dans le cadre de la filière \textbf{INDIA — Ingénierie Numérique, Data
Science et Intelligence Artificielle}.

Cette formation s'inscrit dans le domaine des technologies numériques et met en œuvre différentes
compétences liées au développement informatique, à l'exploitation des données, à l'intelligence
artificielle et à la conception de solutions numériques.

Le projet Campus Navigator constitue une application concrète de plusieurs de ces domaines. Sa
réalisation nécessite notamment de manipuler des données structurées, de concevoir une
architecture logicielle, de développer une interface web, de mettre en œuvre des algorithmes de
recherche de chemins, d'exploiter les capteurs d'un smartphone et d'intégrer des techniques de
traitement du langage naturel.

\subsection{Cadre académique du projet}

La réalisation de Campus Navigator s'inscrit dans le cadre du parcours d'ingénieur en
\textbf{Ingénierie Numérique, Data Science et Intelligence Artificielle}. Elle représente une
occasion de mettre en pratique les connaissances acquises au cours de la formation à travers un
projet répondant à une problématique concrète.

Le projet mobilise plusieurs compétences complémentaires. La gestion et la préparation des
données sont nécessaires pour construire une représentation cohérente du campus. La théorie des
graphes intervient dans la modélisation des connexions entre les emplacements et dans le calcul
des itinéraires. Le développement web permet de construire l'interface et les services de la
solution. Le traitement de signaux issus des capteurs (GPS, accéléromètre) permet de suivre
l'utilisateur pendant son déplacement. Enfin, l'intelligence artificielle est exploitée afin de
faciliter l'interprétation des requêtes formulées en langage naturel et la rédaction des
instructions.

Cette diversité technologique fait du projet un cas d'application transversal permettant de
mettre en relation plusieurs domaines étudiés dans le cadre de la formation.

\subsection{Encadrement du projet}

Le projet a été réalisé sous l'encadrement de \textbf{Madame Alami Kamouri Sophia}.

L'encadrement a permis d'assurer le suivi des différentes étapes du projet, depuis l'étude
initiale des besoins et la préparation des données jusqu'au développement, aux tests et à la
finalisation de la solution.

Un rapport d'avancement a notamment été transmis à l'encadrante le \textbf{3 août}, permettant de
présenter l'état d'avancement du développement et les premières fonctionnalités mises en place.

\section{Contexte général du projet}

\subsection{Contexte de la navigation au sein du campus}

Dans un environnement universitaire, les utilisateurs sont régulièrement amenés à rechercher des
emplacements spécifiques tels que des salles de cours, des laboratoires, des bureaux
administratifs ou différents services. Cette recherche peut être relativement simple pour une
personne connaissant bien l'établissement, mais elle devient plus complexe pour les nouveaux
étudiants, les visiteurs ou toute personne qui ne maîtrise pas l'organisation spatiale du campus.

La difficulté ne consiste pas uniquement à connaître l'emplacement d'une destination. Il est
également nécessaire de déterminer le chemin permettant de l'atteindre depuis une position
donnée, puis de le suivre sans se perdre. Ce chemin peut être composé de plusieurs segments et
nécessiter le passage par différents couloirs, escaliers, bâtiments et allées extérieures.

Afin de répondre à cette problématique, Campus Navigator repose sur une représentation numérique
de l'environnement. Les différents emplacements sont modélisés sous forme de nœuds et leurs
connexions sous forme d'arêtes pondérées par leur longueur. Cette représentation transforme
l'environnement physique du campus en une structure exploitable par des algorithmes de recherche
de chemins.

Parallèlement à cette représentation logique, le projet utilise des plans numériques des
bâtiments. Les plans SVG permettent de représenter visuellement les différents niveaux et de
superposer le parcours calculé, ainsi que la position de l'utilisateur, sur la carte
correspondante.

\subsection{Difficultés rencontrées par les étudiants et visiteurs}

La première difficulté concerne l'identification de la destination. Un utilisateur peut connaître
le nom d'une salle ou d'un service sans connaître son identifiant interne, son niveau ou sa
position exacte sur le plan ; il peut aussi le désigner par un nom d'usage ou avec une faute
d'orthographe.

La deuxième difficulté concerne la détermination de l'itinéraire. Plusieurs chemins peuvent
permettre de rejoindre une destination, et l'utilisateur doit pouvoir identifier le plus court à
partir de sa position de départ.

La présence de plusieurs niveaux et de plusieurs bâtiments constitue également une contrainte
importante. La navigation doit prendre en compte les escaliers, les entrées et les sorties afin
de proposer un parcours cohérent.

Enfin, l'information fournie doit être présentée de manière compréhensible et rester utile
pendant la marche. Une liste de coordonnées ou une représentation technique du graphe ne
constitue pas une information directement exploitable. Le système doit donc transformer le
résultat du calcul en une représentation visuelle, en instructions de déplacement claires et en
un suivi de la progression de l'utilisateur.

\subsection{Besoin d'une solution numérique}

Ces différentes difficultés mettent en évidence l'intérêt d'une solution numérique capable de
centraliser les informations relatives aux lieux et de fournir automatiquement un itinéraire.

Campus Navigator répond à ce besoin en proposant une interaction basée sur une requête formulée en
langage naturel. L'utilisateur indique le lieu qu'il souhaite rejoindre, puis le système identifie
la destination correspondante parmi les emplacements enregistrés.

Une fois la destination identifiée, le système utilise le graphe du campus pour calculer un
itinéraire depuis la position de départ. Le résultat est affiché sur le plan correspondant,
accompagné d'instructions textuelles, puis mis à jour en continu au fur et à mesure que
l'utilisateur avance.

L'approche proposée réunit ainsi dans une même solution la recherche d'un emplacement, le calcul
d'un itinéraire, sa représentation visuelle et le guidage pendant le déplacement.

\section{Problématique}

La problématique principale du projet peut être formulée de la manière suivante :

\begin{quote}
\textbf{Comment concevoir une solution numérique permettant à un utilisateur de localiser
facilement un emplacement au sein d'un campus universitaire, d'obtenir un itinéraire clair à
partir de sa position de départ et d'être guidé jusqu'à destination ?}
\end{quote}

Cette problématique comprend plusieurs aspects complémentaires.

Tout d'abord, le système doit être capable d'interpréter la demande de l'utilisateur. Celui-ci
peut exprimer sa destination sous une forme naturelle, sans connaître la structure interne des
données. Le système doit donc établir une correspondance entre la demande formulée et
l'emplacement correspondant dans la base de données, de manière robuste.

Ensuite, l'environnement du campus doit être représenté sous une forme permettant le calcul
automatique des déplacements. Les emplacements sont considérés comme des nœuds tandis que les
connexions possibles entre ces emplacements sont représentées par des arêtes.

À partir de cette représentation, le système doit déterminer un chemin entre le point de départ
et la destination, puis l'associer aux plans correspondants afin d'en fournir une représentation
visuelle.

Le système doit également estimer la position de l'utilisateur pendant son déplacement, sans
infrastructure de localisation dédiée, et réagir lorsqu'il s'écarte du trajet.

Enfin, l'ensemble du processus doit être présenté à travers une interface simple et intuitive,
accessible depuis un smartphone quel que soit le réseau utilisé.

La problématique du projet se situe ainsi à l'intersection de la \textbf{modélisation spatiale},
de la \textbf{théorie des graphes}, du \textbf{développement web}, de la
\textbf{géolocalisation} et du \textbf{traitement automatique du langage}.

\section{Objectifs du projet}

\subsection{Objectif général}

L'objectif général de \textbf{Campus Navigator} est de développer une solution numérique
permettant d'assister les utilisateurs dans leurs déplacements au sein de l'ENSAM de Rabat, en
leur permettant de rechercher un emplacement, d'obtenir un itinéraire adapté et d'être guidés en
temps réel jusqu'à destination.

La solution cherche notamment à simplifier l'interaction avec le système en permettant à
l'utilisateur de formuler sa demande en langage naturel, puis en transformant cette demande en
une destination exploitable par le moteur de navigation.

\subsection{Objectifs spécifiques}

Pour atteindre cet objectif général, plusieurs objectifs spécifiques ont été définis.

Le premier objectif consiste à \textbf{recenser et structurer les points d'intérêt du campus}.
Au cours du projet, \textbf{68 points d'intérêt} ont été recensés (34 salles de cours,
17 laboratoires, 6 sanitaires, 4 services, 2 bureaux, 2 salles de prière, 2 espaces d'étude et
1 espace événementiel).

Le deuxième objectif consiste à \textbf{préparer et nettoyer les données} nécessaires à la
représentation de l'environnement, organisées autour des niveaux, des emplacements, des nœuds,
des connexions et des types de nœuds.

Le troisième objectif est de \textbf{modéliser l'environnement sous forme d'un graphe pondéré}
afin de permettre le calcul automatique des itinéraires.

Le quatrième objectif consiste à \textbf{développer une API backend} traitant les demandes de
navigation et communiquant avec l'interface utilisateur.

Le cinquième objectif est de \textbf{développer une interface web interactive} permettant de
rechercher une destination, de visualiser le plan et de consulter les indications de
déplacement.

Le sixième objectif consiste à \textbf{intégrer les plans réels du campus} et à assurer leur
calibration afin que le parcours calculé soit correctement positionné sur la représentation
graphique.

Le septième objectif consiste à \textbf{guider l'utilisateur en temps réel}, en estimant sa
position (GPS à l'extérieur, podomètre à l'intérieur) et en recalculant l'itinéraire lorsqu'il
s'en écarte.

Enfin, le projet vise à \textbf{rendre la solution accessible et agréable à utiliser} : support
bilingue, adaptation aux appareils mobiles, utilisation du langage naturel et
\textbf{déploiement sur une adresse publique sécurisée (HTTPS)}.

\section{Présentation de la solution proposée}

\subsection{Principe général de Campus Navigator}

Campus Navigator est une solution de navigation conçue pour faciliter la recherche
d'emplacements au sein de l'ENSAM de Rabat.

L'utilisateur scanne un QR Code affiché sur l'un des écrans du campus, ce qui lui fournit son
point de départ, puis indique dans une interface de chat la destination qu'il souhaite
rejoindre. La requête est transmise au serveur afin d'identifier l'emplacement correspondant dans
les données du campus.

Une fois la destination déterminée, le système recherche le plus court chemin entre le point de
départ et la destination dans le graphe du campus, puis rédige une instruction de marche pour
chaque point de passage. Le parcours est ensuite affiché sur le plan du niveau concerné ; la
carte suit l'utilisateur, efface progressivement la portion déjà parcourue et change
automatiquement de plan lorsque le trajet passe à un autre étage.

Le fonctionnement général peut donc être représenté par la chaîne suivante :

\begin{center}
\textbf{Scan du QR Code $\rightarrow$ Requête utilisateur $\rightarrow$ Identification de la
destination $\rightarrow$ Calcul de l'itinéraire $\rightarrow$ Génération des instructions
$\rightarrow$ Guidage en temps réel sur le plan}
\end{center}

\subsection{Architecture fonctionnelle}

La solution développée repose sur plusieurs composants complémentaires.

Les \textbf{données du campus} constituent la première composante. Elles regroupent les
informations relatives aux niveaux, aux emplacements, aux nœuds, aux connexions et aux types de
nœuds (classeur Excel), complétées par les coordonnées GPS des nœuds extérieurs (fichier JSON).

Le \textbf{backend FastAPI} constitue la deuxième composante. Il fournit l'endpoint
\texttt{/api/navigate} destiné au traitement des demandes de navigation, l'endpoint
\texttt{/api/floorplan} qui renvoie le plan d'un niveau, l'endpoint \texttt{/api/geo} qui
renvoie le réseau extérieur géolocalisé, et l'endpoint \texttt{/api/health} utilisé par la
plateforme d'hébergement. En production, il sert également l'application web compilée.

La troisième composante correspond au \textbf{graphe de navigation}, construit avec NetworkX, qui
représente les relations entre les points du campus et permet de calculer les plus courts
chemins.

La quatrième composante est constituée des \textbf{plans SVG}, qui fournissent une
représentation graphique des espaces et sur lesquels sont tracés le trajet et la position de
l'utilisateur.

Le \textbf{frontend React} fournit l'interface avec laquelle l'utilisateur interagit : une
interface de type chat permet de formuler les demandes, et un module de navigation en temps réel
affiche l'itinéraire et suit l'utilisateur à l'aide des capteurs du téléphone.

Une composante basée sur un \textbf{modèle de langage via Groq} est utilisée pour interpréter
les requêtes et rédiger les instructions de déplacement. Elle est doublée de mécanismes de repli
locaux, de sorte que l'application reste fonctionnelle sans elle.

\subsection{Utilisateurs concernés}

Campus Navigator est principalement destiné aux personnes amenées à se déplacer au sein de
l'établissement.

Les \textbf{étudiants} constituent une catégorie importante d'utilisateurs, notamment lorsqu'ils
doivent rejoindre une salle, un laboratoire ou un service qu'ils connaissent peu.

La solution peut également être utilisée par les \textbf{enseignants et le personnel
administratif}, notamment pour rechercher des espaces qu'ils fréquentent occasionnellement.

Les \textbf{visiteurs} peuvent également en bénéficier lors d'événements ou de rencontres
organisés au sein de l'établissement.

Enfin, la solution peut constituer un outil exploitable par l'\textbf{administration} pour
centraliser et gérer les informations relatives aux lieux du campus.

\subsection{Valeur apportée par la solution}

L'intérêt principal de Campus Navigator réside dans la combinaison de plusieurs fonctionnalités
au sein d'une même application.

La solution ne se limite pas à fournir une représentation statique du campus : elle calcule
automatiquement un parcours à partir d'une position de départ, puis accompagne l'utilisateur
pendant sa marche.

L'utilisation d'une interface conversationnelle rend la recherche plus naturelle : l'utilisateur
exprime sa demande sans connaître les identifiants techniques utilisés par le système.

La représentation du trajet sur les plans SVG, combinée au suivi de la position, apporte une
dimension visuelle importante et facilite la compréhension du parcours.

Enfin, aucune installation n'est requise : l'application s'ouvre dans le navigateur du téléphone
par simple scan d'un QR Code, et ne nécessite aucune infrastructure de localisation (balises,
bornes) sur le campus.

\section{Planification du projet}

\subsection{Méthodologie de réalisation}

La réalisation du projet Campus Navigator s'est déroulée du \textbf{8 juillet au 30 août} selon
une démarche progressive, puis a été prolongée fin octobre par une phase d'extension. Le travail
a été organisé en plusieurs phases allant de l'étude des besoins au déploiement de la solution.

La première phase a été consacrée à l'étude des besoins et à la collecte des données. Cette étape
a permis d'identifier les utilisateurs cibles, de recenser les points d'intérêt du campus et de
préparer les données nécessaires à la navigation. Au total, \textbf{68 points d'intérêt} ont été
recensés. La base de données Excel a ensuite été construite et nettoyée autour de plusieurs
ensembles d'informations : niveaux, emplacements, nœuds, arêtes et types de nœuds. Un contrôle de
la qualité des données a également été réalisé à l'aide de Power BI.

La deuxième phase, réalisée du \textbf{21 au 27 juillet}, a porté sur la modélisation du graphe de
navigation avec NetworkX. Le graphe comporte aujourd'hui \textbf{145 nœuds} et \textbf{174
arêtes} saisies (172 connexions distinctes une fois les doublons fusionnés). Une analyse de sa
structure a permis d'identifier les nœuds isolés et de procéder aux corrections nécessaires afin
d'améliorer la cohérence du réseau.

La troisième phase, du \textbf{28 juillet au 3 août}, a été consacrée au développement de
l'application. Le backend a été développé avec FastAPI et l'endpoint \texttt{/api/navigate} a
été mis en place. En parallèle, une interface de chat a été développée avec React, et les QR Codes
des écrans du campus ont été générés. Une première intégration d'un modèle de langage, via l'API
Gemini, a permis d'exploiter le langage naturel ; c'est cette version qui figure dans le rapport
d'avancement. Les limites de quota de cette API ayant rendu son usage impossible, l'application a
ensuite été migrée vers l'API Groq (voir section~\ref{sec:llm}). Un rapport d'avancement a été
envoyé à l'encadrante le \textbf{3 août}.

La quatrième phase, du \textbf{4 au 10 août}, a porté sur la cartographie et la calibration. Les
cartes simplifiées utilisées initialement ont été remplacées par les plans SVG réels. Un travail
d'alignement précis des coordonnées a été effectué afin de garantir la correspondance entre le
graphe et la représentation visuelle du campus. Plusieurs problèmes d'affichage ont également été
corrigés, notamment l'absence de visibilité du tracé et la présence de doublons dans les onglets
associés aux étages.

La cinquième phase, du \textbf{11 au 17 août}, a été consacrée à l'amélioration de l'expérience
utilisateur : navigation tactile (déplacement du plan, zoom par pincement), meilleure visibilité
du tracé et ajout de marqueurs pour la position et la destination.

La sixième phase, du \textbf{18 au 24 août}, a regroupé le support bilingue (\textbf{français et
anglais}), les tests et la finalisation de l'application. Des tests utilisateurs ont permis
d'identifier les derniers problèmes fonctionnels et ergonomiques, corrigés avant la mise en
service de la solution en réseau local.

La période du \textbf{25 au 30 août} a été consacrée à la rédaction d'une première version du
rapport et à la préparation de la présentation.

Enfin, une phase d'extension, réalisée \textbf{fin octobre 2026}, a porté sur le guidage en temps
réel et la mise en ligne : refonte de l'interface pour mobile, suivi de la position par
accéléromètre puis par GPS sur les trajets extérieurs, recalcul automatique de l'itinéraire,
amélioration de la rédaction des instructions, et déploiement de l'application sur une adresse
HTTPS publique avec régénération des QR Codes.

\subsection{Synthèse des étapes réalisées}

Le déroulement du projet peut être résumé dans le tableau suivant :

\begin{table}[h]
\centering
\caption{Synthèse des étapes réalisées}
\label{tab:etapes}
\small
\begin{tabularx}{\linewidth}{@{}p{2.6cm}p{3.2cm}X@{}}
\toprule
\textbf{Période} & \textbf{Phase} & \textbf{Principales réalisations}\\
\midrule
8 -- 20 juillet & Étude des besoins et collecte & 68 points d'intérêt, préparation et nettoyage des données, contrôle Power BI\\
21 -- 27 juillet & Modélisation du graphe & 145 nœuds, 174 arêtes, correction des nœuds isolés\\
28 juillet -- 3 août & Développement & FastAPI, API de navigation, React, QR Codes, première intégration LLM\\
4 -- 10 août & Cartographie et calibration & Plans SVG réels, alignement des coordonnées, corrections d'affichage\\
11 -- 17 août & Amélioration UX & Support tactile, pan/zoom, visibilité du trajet, marqueurs\\
18 -- 24 août & Bilinguisme, tests et finalisation & Français et anglais, tests utilisateurs, corrections\\
25 -- 30 août & Documentation & Première rédaction du rapport\\
Fin octobre & Navigation temps réel et déploiement & GPS extérieur, podomètre, recalcul d'itinéraire, Docker, hébergement HTTPS, nouveaux QR Codes\\
\bottomrule
\end{tabularx}
\end{table}

\subsection{Diagramme de Gantt}

Afin de représenter visuellement la répartition temporelle des activités, un diagramme de Gantt
a été élaboré pour la période de stage (Figure~\ref{fig:gantt}). La phase d'extension de fin
octobre, réalisée après la période de stage, n'y figure pas.

\begin{figure}[h]
\centering
\resizebox{\linewidth}{!}{%
\begin{ganttchart}[
    time slot format=isodate,
    x unit=0.28cm, y unit chart=0.7cm,
    hgrid, vgrid={*6{draw=none},dotted},
    title height=1,
    bar/.append style={fill=blue!60!black},
    bar height=0.5,
    milestone/.append style={fill=black},
    bar label font=\small, milestone label font=\small\itshape,
  ]{2026-07-08}{2026-08-30}
  \gantttitle{Juillet 2026}{24} \gantttitle{Août 2026}{30} \\
  \ganttbar{Étude des besoins \& collecte}{2026-07-08}{2026-07-20} \\
  \ganttbar{Modélisation du graphe}{2026-07-21}{2026-07-27} \\
  \ganttbar[bar/.append style={fill=orange!85!black}]{Développement}{2026-07-28}{2026-08-03} \\
  \ganttmilestone{Rapport d'avancement}{2026-08-03} \\
  \ganttbar[bar/.append style={fill=orange!85!black}]{Cartographie \& calibration}{2026-08-04}{2026-08-10} \\
  \ganttbar[bar/.append style={fill=orange!85!black}]{Amélioration UX}{2026-08-11}{2026-08-17} \\
  \ganttbar[bar/.append style={fill=orange!85!black}]{Bilinguisme, tests \& finalisation}{2026-08-18}{2026-08-24} \\
  \ganttmilestone{Mise en service locale}{2026-08-24} \\
  \ganttbar[bar/.append style={fill=black!80}]{Rédaction du rapport}{2026-08-25}{2026-08-30}
\end{ganttchart}}
\caption{Diagramme de Gantt du projet Campus Navigator}
\label{fig:gantt}
\end{figure}

Le diagramme permet de visualiser l'enchaînement des phases et leur répartition sur la période du
8 juillet au 30 août. Cette organisation progressive a permis de faire évoluer le projet depuis la
préparation des données et la modélisation du campus jusqu'à l'intégration des composantes, aux
tests et à la finalisation de la solution.

\section{Conclusion}

Ce chapitre a permis de présenter le cadre général du projet \textbf{Campus Navigator}, réalisé
au sein de l'\textbf{ENSAM de Rabat} dans le cadre de la filière \textbf{INDIA — Ingénierie
Numérique, Data Science et Intelligence Artificielle}, sous l'encadrement de \textbf{Madame
Alami Kamouri Sophia}.

L'étude du contexte a mis en évidence la difficulté que peut représenter la recherche d'un
emplacement au sein d'un campus comportant plusieurs bâtiments et niveaux, en particulier pour
les utilisateurs qui ne connaissent pas l'organisation spatiale de l'établissement.

Pour répondre à cette problématique, Campus Navigator combine une représentation structurée des
données du campus, une modélisation sous forme de graphe, un calcul d'itinéraires, des plans SVG,
une interface web interactive et un guidage en temps réel exploitant les capteurs du téléphone.
L'intégration d'un modèle de langage facilite l'interaction à travers des requêtes en langage
naturel.

Le projet a été réalisé progressivement entre le \textbf{8 juillet et le 30 août}, puis étendu
fin octobre, en passant par l'étude des besoins, la préparation des données, la modélisation du
graphe, le développement, la calibration des plans, l'amélioration de l'expérience utilisateur,
le support bilingue, les tests, le guidage en temps réel et le déploiement.

Le chapitre suivant est consacré à l'\textbf{état de l'art}. Il étudie les principales approches,
technologies et méthodes utilisées dans le domaine de la navigation intérieure et situe les choix
technologiques effectués dans le cadre de Campus Navigator.

% =====================================================================
% Chapitre 2
% =====================================================================
\chapter{État de l'art}

\section{Introduction}

Le développement d'une solution de navigation destinée à un environnement universitaire nécessite
la prise en compte de plusieurs domaines technologiques complémentaires. Contrairement à la
navigation extérieure, qui repose largement sur des systèmes de positionnement satellitaire et
sur des cartographies standardisées, la navigation intérieure doit s'adapter à des espaces
complexes, composés de plusieurs bâtiments, niveaux, couloirs, salles, escaliers et points
d'accès. Dans ce contexte, la représentation de l'environnement, la détermination de la position
de l'utilisateur et le calcul d'un itinéraire adapté constituent des problématiques essentielles.

Les technologies de localisation telles que le Wi-Fi, le Bluetooth Low Energy, l'Ultra-Wideband,
les capteurs inertiels des smartphones ou encore les QR Codes peuvent être utilisées pour
déterminer ou faciliter la détermination de la position d'un utilisateur. Parallèlement, les
techniques de cartographie numérique permettent de représenter les espaces intérieurs, tandis que
les algorithmes issus de la théorie des graphes permettent de calculer des itinéraires.

Plus récemment, l'intelligence artificielle et les modèles de langage ont ouvert de nouvelles
possibilités pour améliorer l'interaction entre l'utilisateur et les systèmes de navigation. Au
lieu de rechercher manuellement une destination dans une liste, l'utilisateur peut formuler sa
demande en langage naturel ; le système doit alors interpréter cette requête, identifier la
destination et produire un itinéraire compréhensible.

Dans ce chapitre, l'état de l'art est présenté à travers les principales technologies et méthodes
susceptibles d'être utilisées dans une solution de navigation intérieure : caractéristiques de la
navigation indoor, technologies de localisation, représentation cartographique, calcul
d'itinéraires, intelligence artificielle et langage naturel, technologies de développement et de
déploiement web, et stockage des données. Une comparaison des approches existantes permet enfin
de positionner Campus Navigator.

\section{Navigation intérieure}

\subsection{Définition de la navigation intérieure}

La navigation intérieure, ou \emph{indoor navigation}, désigne l'ensemble des techniques
permettant à un utilisateur de se déplacer efficacement à l'intérieur d'un bâtiment ou d'un
complexe de bâtiments. Elle comprend généralement trois fonctions complémentaires : la
localisation de l'utilisateur, la détermination de la destination et le calcul d'un itinéraire
entre la position de départ et cette destination ; une quatrième fonction, le suivi pendant le
déplacement, permet de passer d'un simple calcul d'itinéraire à un véritable guidage.

Dans un environnement universitaire, la navigation intérieure peut concerner des salles de cours,
des laboratoires, des bureaux administratifs, des amphithéâtres, des bibliothèques ou des espaces
communs. La diversité de ces espaces rend la recherche d'un emplacement plus complexe,
particulièrement pour les nouveaux étudiants ou les visiteurs.

La navigation intérieure ne se limite donc pas à afficher une carte. Une solution complète doit
représenter la structure du bâtiment, identifier les lieux accessibles et déterminer un chemin
cohérent entre deux points.

\subsection{Particularités et contraintes}

La principale difficulté de la navigation intérieure réside dans l'absence d'une référence de
positionnement universelle comparable au GPS utilisé en extérieur. Les signaux satellitaires sont
fortement atténués ou indisponibles à l'intérieur des bâtiments. Il devient alors nécessaire
d'utiliser d'autres technologies ou mécanismes pour déterminer la position de l'utilisateur.

La structure physique des bâtiments constitue également une contrainte importante. Un bâtiment
peut comporter plusieurs étages reliés par des escaliers, des ascenseurs ou des passages
spécifiques. Le système doit tenir compte de la connectivité entre les niveaux et éviter de
proposer des chemins qui traversent des murs ou des zones non accessibles.

Enfin, l'information fournie doit être compréhensible. Une représentation graphique précise ne
suffit pas toujours à garantir une bonne expérience de navigation ; les indications doivent
également être présentées clairement, ce qui explique l'intérêt croissant des interfaces
interactives et du langage naturel.

\section{Technologies de localisation intérieure}

La localisation constitue l'un des composants fondamentaux d'un système de navigation. Elle
permet de déterminer le point à partir duquel l'itinéraire doit être calculé, puis de suivre
l'utilisateur. Plusieurs technologies peuvent être utilisées, chacune présentant un compromis
différent entre précision, coût, infrastructure et facilité de déploiement.

\subsection{GPS et GNSS}

Le GPS, intégré plus largement dans la famille des systèmes GNSS, constitue une technologie de
référence pour la localisation extérieure. Il permet d'estimer la position d'un appareil à partir
des signaux provenant de satellites. Dans un navigateur web, il est accessible via l'API de
géolocalisation (\texttt{navigator.geolocation}), qui fournit la latitude, la longitude et une
estimation de la précision ; cette API n'est disponible que sur des pages servies en HTTPS.

Son utilisation à l'intérieur des bâtiments est limitée par l'atténuation des signaux : il ne
permet généralement pas de déterminer la salle ou le couloir dans lequel se trouve l'utilisateur.

Le GPS reste néanmoins pertinent sur les parties extérieures d'un campus (allées, cours, accès
entre bâtiments), où une précision de quelques mètres suffit à suivre l'utilisateur. C'est
l'usage qui en est fait dans Campus Navigator.

\subsection{Wi-Fi}

Les réseaux Wi-Fi déjà présents dans de nombreux bâtiments peuvent être exploités pour la
localisation intérieure. La position peut être estimée à partir de la puissance des signaux reçus
ou grâce à une phase de calibration appelée \emph{fingerprinting}.

L'un des principaux avantages du Wi-Fi réside dans la réutilisation d'une infrastructure
existante. Cependant, la précision dépend fortement de l'environnement et de la disposition des
points d'accès, et les modifications de l'infrastructure peuvent nécessiter une nouvelle
calibration. En outre, les navigateurs web n'exposent pas la puissance des signaux Wi-Fi, ce qui
réserve cette approche aux applications natives.

\subsection{Bluetooth Low Energy}

Le Bluetooth Low Energy, notamment à travers les balises (\emph{beacons}), constitue une autre
approche couramment étudiée. Des balises fixes peuvent être installées dans différentes zones du
bâtiment et leur signal détecté par le smartphone de l'utilisateur.

Cette technologie présente l'avantage d'une consommation énergétique faible et d'une
infrastructure pouvant être déployée progressivement. Toutefois, la précision dépend de la
configuration des balises et des conditions physiques de l'environnement, et l'installation
représente un coût matériel.

\subsection{Ultra-Wideband}

L'Ultra-Wideband (UWB) utilise des signaux radio permettant d'obtenir des mesures de distance
très précises. Cette technologie est intéressante lorsque la précision de localisation constitue
une exigence importante.

Son principal obstacle dans le cadre d'un projet universitaire réside dans les besoins matériels :
son déploiement nécessite des équipements compatibles et une infrastructure adaptée.

\subsection{Capteurs inertiels et podométrie}

Les smartphones intègrent des capteurs inertiels (accéléromètre, gyroscope, magnétomètre),
accessibles depuis le navigateur via l'événement \texttt{devicemotion}. La navigation à l'estime
pédestre (\emph{Pedestrian Dead Reckoning}, PDR) exploite ces capteurs pour détecter les pas de
l'utilisateur et estimer la distance parcourue, sans aucune infrastructure.

Cette approche présente l'avantage d'être disponible sur pratiquement tous les téléphones. Sa
limite principale est la dérive : l'erreur s'accumule avec la distance (longueur de pas variable,
pas manqués ou détectés à tort). Elle est donc particulièrement adaptée lorsqu'elle est combinée à
une contrainte forte, comme le fait de savoir que l'utilisateur suit un itinéraire connu, et à des
points de recalage (GPS, QR Codes). Sur iOS, l'accès à ces capteurs doit être autorisé
explicitement par l'utilisateur, à la suite d'une interaction (un appui).

\subsection{QR Codes et localisation assistée}

Les QR Codes constituent une approche différente. Plutôt que d'estimer automatiquement la position
de l'utilisateur, le système lui demande de scanner un code placé à un emplacement connu. Le scan
associe alors l'utilisateur à un nœud déterminé de la carte.

Cette méthode est particulièrement intéressante pour un prototype : elle est peu coûteuse, facile
à déployer et fournit une position initiale exacte. En revanche, elle nécessite une action de
l'utilisateur et ne fournit pas, à elle seule, un suivi continu ; elle peut cependant servir de
point de recalage si des codes sont placés le long des trajets (par exemple aux escaliers).

Un QR Code est un code-barres bidimensionnel normalisé (ISO/IEC 18004), capable d'encoder une
chaîne de caractères telle qu'une URL. Dans Campus Navigator, chaque code encode l'adresse
publique de l'application paramétrée avec l'identifiant du nœud de départ
(\texttt{https://<adresse>/?start=N114}) ; sa génération est automatisée par un script Python
(bibliothèque \texttt{qrcode}).

\subsection{Comparaison des technologies de localisation}

Le choix d'une technologie dépend du niveau de précision recherché, de l'infrastructure disponible
et du contexte d'utilisation.

\begin{table}[h]
\centering
\caption{Comparaison synthétique des technologies de localisation intérieure}
\label{tab:loc}
\small
\begin{tabularx}{\linewidth}{@{}lXX@{}}
\toprule
\textbf{Technologie} & \textbf{Principal intérêt} & \textbf{Principale limite}\\
\midrule
\textbf{GPS/GNSS} & Disponible sur tout smartphone, efficace en extérieur & Faible efficacité à l'intérieur\\
\textbf{Wi-Fi} & Utilise une infrastructure existante & Précision variable, inaccessible depuis un navigateur\\
\textbf{BLE} & Faible consommation et déploiement flexible & Installation de balises nécessaire\\
\textbf{UWB} & Localisation très précise & Coût et matériel spécialisé\\
\textbf{Capteurs inertiels (PDR)} & Aucune infrastructure, disponible partout & Dérive avec la distance\\
\textbf{QR Code} & Simple, peu coûteux, position exacte & Nécessite une action de l'utilisateur\\
\bottomrule
\end{tabularx}
\end{table}

Dans le contexte de Campus Navigator, une approche hybride a été retenue : le QR Code fournit une
position de départ exacte, le GPS suit l'utilisateur sur les tronçons extérieurs, et la podométrie
prend le relais à l'intérieur, la progression étant contrainte à l'itinéraire calculé. Cette
combinaison ne nécessite aucune infrastructure supplémentaire sur le campus.

\section{Cartographie et représentation spatiale}

\subsection{Cartographie numérique}

La cartographie numérique consiste à représenter un environnement réel sous une forme exploitable
par un système informatique. Dans le contexte de la navigation intérieure, cette représentation
doit permettre d'identifier les espaces, les passages et les connexions entre les différentes
zones du bâtiment.

Une carte intérieure doit notamment représenter les éléments nécessaires au déplacement : salles,
couloirs, escaliers, ascenseurs, entrées et autres points importants. Elle peut également contenir
des informations sémantiques associant une position à une fonction (« laboratoire », « bureau »,
« amphithéâtre »).

\subsection{Cartographie vectorielle}

Les données vectorielles représentent les éléments géographiques à l'aide de points, de lignes et
de polygones. Cette approche est particulièrement adaptée aux plans de bâtiments car elle permet de
représenter précisément les contours et les connexions entre les éléments.

Contrairement à une image matricielle, une représentation vectorielle peut être agrandie sans
perte de qualité — propriété essentielle lorsque l'utilisateur zoome sur le plan — et peut être
manipulée directement par une application web.

\subsection{SVG et Inkscape}

Le format SVG (\emph{Scalable Vector Graphics}) est particulièrement adapté à l'affichage de plans
interactifs dans une application web. Un plan SVG peut être affiché comme fond de carte, et une
seconde couche SVG partageant le même repère (\texttt{viewBox}) peut y superposer l'itinéraire, la
position de l'utilisateur et des marqueurs.

Dans Campus Navigator, les plans des quatre niveaux ont été conçus avec l'éditeur vectoriel
Inkscape, puis intégrés dans l'interface web. Cette approche permet de conserver une
représentation visuelle légère et personnalisable sans introduire la complexité d'un SIG complet.

\begin{table}[h]
\centering
\caption{Principales approches de représentation cartographique}
\label{tab:carto}
\small
\begin{tabularx}{\linewidth}{@{}lXX@{}}
\toprule
\textbf{Approche} & \textbf{Caractéristique} & \textbf{Usage possible}\\
\midrule
\textbf{Image raster} & Image fixe composée de pixels & Plan simple\\
\textbf{SVG} & Format vectoriel interactif & Plan intérieur interactif\\
\textbf{GeoJSON} & Données géographiques structurées & Données spatiales\\
\textbf{Vector tiles} & Données vectorielles distribuées par tuiles & Cartographie à grande échelle\\
\bottomrule
\end{tabularx}
\end{table}

Pour un environnement universitaire de taille limitée, le SVG présente un compromis intéressant
entre précision graphique, simplicité de manipulation et intégration dans une application web.

\section{Modélisation et calcul d'itinéraires}

\subsection{Théorie des graphes appliquée à la navigation}

Le calcul d'un itinéraire peut être modélisé comme un problème de recherche de chemin dans un
graphe $G = (V, E)$. Les lieux ou points de passage sont représentés par des sommets $v \in V$ et
les connexions possibles par des arêtes $e \in E$.

Dans un bâtiment, un sommet peut correspondre à une salle, une intersection de couloir, une entrée
ou un escalier. Une arête représente une possibilité de déplacement entre deux sommets et est
associée à un poids $w(e) \geq 0$ représentant une distance ou un coût de déplacement.

Cette modélisation transforme le problème physique de navigation en un problème algorithmique de
recherche de chemin.

\subsection{Algorithme de Dijkstra}

L'algorithme de Dijkstra est un algorithme classique de recherche des plus courts chemins dans un
graphe pondéré à poids non négatifs. Il calcule le chemin $P = (v_0, \ldots, v_k)$ minimisant le
coût total $\sum_{i=0}^{k-1} w(v_i, v_{i+1})$ entre un point de départ et une destination.

Si les poids représentent les distances entre les différents points, le résultat correspond au
chemin de distance minimale. Son principal avantage est sa simplicité et sa robustesse ; sur un
graphe de quelques centaines de sommets comme celui d'un campus, son temps d'exécution est
négligeable.

\subsection{Algorithme A*}

L'algorithme A* combine le coût déjà parcouru avec une estimation (heuristique) du coût restant
jusqu'à la destination. Grâce à cette heuristique, il explore une partie plus ciblée du graphe que
Dijkstra lorsque l'heuristique est adaptée. Il est largement utilisé dans les systèmes de
navigation et les environnements interactifs de grande taille.

\subsection{Bibliothèques et moteurs de routage}

Plusieurs bibliothèques et moteurs spécialisés peuvent être utilisés pour implémenter la recherche
d'itinéraires. NetworkX fournit des structures de graphes et plusieurs algorithmes de recherche de
chemin directement utilisables en Python.

Pour des systèmes plus importants ou nécessitant une infrastructure géospatiale avancée, des
solutions telles que pgRouting, GraphHopper ou OSRM peuvent être envisagées, selon la taille du
graphe, le volume de requêtes et le niveau d'intégration avec les données géographiques.

\begin{table}[h]
\centering
\caption{Comparaison des principales approches de routage}
\label{tab:routage}
\small
\begin{tabularx}{\linewidth}{@{}lXX@{}}
\toprule
\textbf{Approche} & \textbf{Principe} & \textbf{Utilisation}\\
\midrule
\textbf{Dijkstra} & Recherche du plus court chemin & Graphes pondérés\\
\textbf{A*} & Recherche guidée par une heuristique & Navigation à grande échelle\\
\textbf{NetworkX} & Bibliothèque Python de graphes & Prototypage et applications Python\\
\textbf{pgRouting} & Routage associé à PostgreSQL/PostGIS & Données spatiales importantes\\
\textbf{GraphHopper / OSRM} & Moteurs de routage spécialisés & Services de navigation\\
\bottomrule
\end{tabularx}
\end{table}

Dans Campus Navigator, la représentation du campus sous forme de graphe dissocie clairement les
données cartographiques de la logique de calcul d'itinéraire : le système identifie une
destination puis recherche le chemin correspondant dans le graphe.

\section{Intelligence artificielle et interaction en langage naturel}

\subsection{Intelligence artificielle dans les systèmes de navigation}

L'intelligence artificielle peut intervenir à différentes étapes d'un système de navigation :
analyse des préférences de l'utilisateur, interprétation d'informations textuelles ou génération
des instructions de déplacement. Dans le cadre d'une navigation intérieure, elle est
particulièrement intéressante lorsque l'utilisateur interagit avec le système en langage naturel.

\subsection{Modèles de langage}

Les grands modèles de langage (LLM) sont capables d'analyser des requêtes formulées naturellement
et d'en extraire des informations pertinentes. Un utilisateur peut exprimer une demande sous une
forme différente de celle utilisée dans la base de données ; le modèle peut alors identifier
l'intention et faire correspondre la requête à une destination connue. Il peut également
reformuler une suite de points de passage en instructions de marche naturelles.

L'utilisation d'un LLM ne remplace cependant pas le moteur de routage. Le rôle du modèle est limité
à l'interprétation de la requête et à la rédaction, tandis que le calcul du chemin reste effectué
par un algorithme déterministe. Les sorties du modèle doivent par ailleurs être contrôlées, car il
peut produire des réponses mal formées ou inventer des informations.

\subsection{Recherche sémantique et correspondance approximative}

Lorsque les noms des lieux sont nombreux ou que les utilisateurs emploient des formulations
différentes, une recherche sémantique (embeddings, bases vectorielles) peut améliorer la
correspondance entre une requête et les destinations disponibles.

Pour un nombre limité de destinations, des techniques plus simples peuvent suffire : normalisation
du texte (minuscules, suppression des accents et de la ponctuation), recherche des mots-clés d'un
nom de lieu dans la requête, et mesure de similarité entre chaînes de caractères (par exemple le
ratio de la bibliothèque Python \texttt{difflib}) pour tolérer les fautes de frappe. C'est cette
approche qui sert de mécanisme de repli dans Campus Navigator.

\subsection{Services d'inférence LLM}

Des services d'inférence tels que Groq donnent accès à des modèles de langage hébergés sur une
infrastructure externe, ce qui évite d'exécuter localement un modèle volumineux.

L'utilisation d'un service externe introduit des contraintes liées à la connexion Internet, à la
latence, aux limites d'utilisation et à la confidentialité des données transmises. Une
architecture robuste doit donc prévoir un comportement de secours lorsque le service n'est pas
disponible.

Dans Campus Navigator, plusieurs modèles ont été utilisés successivement. Une première version
reposait sur l'API Gemini de Google, abandonnée en raison de limites de quota (voir
section~\ref{sec:llm}), au profit de l'API Groq. Le premier modèle retenu,
\texttt{llama-3.1-8b-instant} (famille Llama~3.1 de Meta, 8 milliards de paramètres), a été
choisi pour sa faible latence. Il a ensuite été remplacé par \texttt{openai/gpt-oss-20b}, puis par
\texttt{openai/gpt-oss-120b}, modèles à poids ouverts publiés par OpenAI en 2025, la version de
plus grande taille produisant des instructions nettement plus fiables. Le modèle est désormais un
simple paramètre de configuration (variable d'environnement \texttt{GROQ\_MODEL}), sans
modification de l'architecture d'appel.

\section{Technologies de développement web}

\subsection{Architecture client-serveur}

Une application de navigation web peut être organisée selon une architecture client-serveur. Le
client correspond à l'interface utilisée pour rechercher une destination, visualiser le plan et
suivre le guidage ; le serveur prend en charge la logique applicative et le traitement des
requêtes. Cette séparation rend les composants indépendants et facilite leur évolution.

\subsection{React}

React est une bibliothèque JavaScript destinée à la construction d'interfaces utilisateur. Son
approche basée sur des composants et des \emph{hooks} permet de décomposer une application en
éléments réutilisables et de gérer un état complexe : historique de conversation, position de
l'utilisateur, caméra de la carte, abonnements aux capteurs.

\subsection{Vite}

Vite est un outil de développement et de construction d'applications web. Il fournit un serveur de
développement rapide et produit une version optimisée de l'application pour la production. Dans
Campus Navigator, il est configuré pour servir l'application en HTTPS pendant le développement et
pour relayer les appels \texttt{/api} vers le backend.

\subsection{FastAPI}

FastAPI est un framework Python destiné au développement d'API web. Il s'appuie sur Pydantic pour
la validation des données reçues et génère automatiquement la documentation des endpoints. Dans
Campus Navigator, il reçoit les demandes de navigation, interroge le graphe, appelle le modèle de
langage et sert également, en production, les fichiers statiques de l'application React.

\subsection{Node.js}

Node.js est un environnement d'exécution JavaScript fonctionnant en dehors du navigateur. Il est
accompagné du gestionnaire de paquets npm et constitue l'outillage sur lequel reposent Vite et
React lors du développement et de la compilation. Dans Campus Navigator, il intervient uniquement
à la construction du frontend ; la génération des QR Codes, initialement réalisée avec Node.js, a
été réécrite en Python afin de lire directement le classeur de données.

\subsection{Conteneurisation et hébergement}

Docker permet d'empaqueter une application et ses dépendances dans une image exécutable à
l'identique sur n'importe quel serveur. Une construction en plusieurs étapes (\emph{multi-stage
build}) permet de compiler le frontend dans une image Node.js, puis de ne conserver que le résultat
dans une image Python légère qui exécute l'API.

Des plateformes d'hébergement telles que Render déploient directement une image Docker à partir
d'un dépôt Git et fournissent automatiquement une adresse HTTPS, condition nécessaire pour
accéder au GPS et aux capteurs de mouvement depuis un navigateur mobile.

\begin{table}[h]
\centering
\caption{Technologies utilisées dans l'architecture web}
\label{tab:web}
\small
\begin{tabularx}{\linewidth}{@{}lXX@{}}
\toprule
\textbf{Technologie} & \textbf{Rôle principal} & \textbf{Contribution au projet}\\
\midrule
\textbf{React} & Interface utilisateur & Chat, carte, guidage en temps réel\\
\textbf{Vite} & Environnement de développement & Développement HTTPS, compilation\\
\textbf{FastAPI} & API backend & Communication, logique serveur, service du site\\
\textbf{NetworkX} & Traitement des graphes & Calcul des itinéraires\\
\textbf{SVG} & Représentation graphique & Affichage des plans\\
\textbf{Groq / LLM} & Traitement du langage & Interprétation des demandes, rédaction des instructions\\
\textbf{API Web (Geolocation, DeviceMotion)} & Capteurs du téléphone & Suivi de la position\\
\textbf{Python (qrcode)} & Script utilitaire & Génération des QR Codes\\
\textbf{Docker / Render} & Conteneurisation, hébergement & Mise en ligne en HTTPS\\
\bottomrule
\end{tabularx}
\end{table}

Cette combinaison sépare les responsabilités : React prend en charge l'expérience utilisateur et
le suivi de la position, FastAPI assure la logique serveur, tandis que les composants de calcul et
d'intelligence artificielle réalisent les traitements spécifiques.

\section{Gestion et stockage des données}

\subsection{Données nécessaires}

Un système de navigation intérieure repose sur plusieurs catégories de données : les lieux, les
connexions entre les lieux et les caractéristiques des niveaux du bâtiment.

Les données relatives aux lieux contiennent un identifiant, un nom, une catégorie et une position
sur le plan. Les données relatives aux connexions permettent de construire le graphe utilisé par le
moteur de routage.

La qualité de ces données est déterminante. Un itinéraire calculé correctement à partir d'un graphe
incomplet ou incorrect peut malgré tout être inutilisable dans le monde réel.

\subsection{Fichiers tabulaires}

Les fichiers tabulaires tels que les classeurs Excel constituent une solution simple pour stocker et
modifier les informations pendant une phase de prototypage. Cette approche est accessible à des
utilisateurs non spécialisés, mais devient moins adaptée lorsque le volume de données augmente ou
que plusieurs personnes doivent modifier les données simultanément.

Pour exploiter ces fichiers dans le backend, la bibliothèque Python \textbf{pandas} fournit la
structure DataFrame. Dans Campus Navigator, pandas charge au démarrage du serveur les feuilles
\emph{Floors}, \emph{Nodes}, \emph{Edges} et \emph{Locations}, tandis que la feuille
\emph{Node\_Types} est lue par le script de génération des QR Codes. Les coordonnées GPS des nœuds
extérieurs sont stockées séparément dans un fichier JSON.

\subsection{Bases de données spatiales}

Les bases de données spatiales telles que PostgreSQL avec l'extension PostGIS permettent de stocker
et de manipuler des données géographiques de manière structurée. Elles deviennent pertinentes
lorsque le projet doit gérer plusieurs bâtiments, des géométries complexes ou des requêtes
spatiales avancées.

\subsection{Bases orientées graphes}

Les bases de données orientées graphes représentent directement les relations entre les éléments.
Elles sont intéressantes pour des environnements complexes dans lesquels les relations entre les
lieux constituent une partie essentielle du modèle de données.

\begin{table}[h]
\centering
\caption{Comparaison des solutions de stockage}
\label{tab:stockage}
\small
\begin{tabularx}{\linewidth}{@{}lXX@{}}
\toprule
\textbf{Solution} & \textbf{Avantage principal} & \textbf{Limite principale}\\
\midrule
\textbf{Excel / fichiers tabulaires} & Simplicité de modification & Scalabilité limitée\\
\textbf{PostgreSQL / PostGIS} & Gestion structurée des données spatiales & Mise en place plus complexe\\
\textbf{Base orientée graphe} & Représentation naturelle des relations & Infrastructure spécialisée\\
\bottomrule
\end{tabularx}
\end{table}

Pour un prototype de taille limitée, une structure tabulaire constitue donc un choix pragmatique.
Une évolution vers une base spécialisée pourrait être envisagée si le nombre de bâtiments, de lieux
et de requêtes augmente.

\subsection{Contrôle de la qualité des données et Power BI}

\textbf{Power BI} est un outil d'analyse et de visualisation de données développé par Microsoft.
Il permet d'importer des sources variées (dont des classeurs Excel), de définir un modèle de données
relationnel entre plusieurs tables, puis de construire des tableaux de bord interactifs. Les
indicateurs et colonnes calculées y sont définis grâce au langage \textbf{DAX} (\emph{Data Analysis
Expressions}).

Bien qu'il soit principalement destiné à l'analyse décisionnelle, Power BI peut être utilisé pour
auditer la qualité d'un jeu de données : détection de doublons, de références manquantes ou de
valeurs incohérentes. Cet usage est pertinent dans un projet de navigation, où la qualité du graphe
conditionne directement la fiabilité des itinéraires. Dans Campus Navigator, il a servi à contrôler
la cohérence du classeur Excel avant et pendant le développement (voir section~\ref{sec:bdd}).

\section{Étude comparative des solutions existantes}

\subsection{Solutions commerciales de navigation intérieure}

Le marché comprend plusieurs plateformes proposant des fonctionnalités de cartographie et de
navigation intérieure. Des solutions telles que Mapbox, HERE, Esri Indoors ou Pointr illustrent
différentes approches permettant d'intégrer des cartes intérieures dans des applications
numériques.

Ces plateformes proposent des fonctionnalités avancées de cartographie, de gestion des bâtiments
ou de localisation, et fournissent des composants déjà développés. Cependant, leur utilisation peut
nécessiter des données cartographiques spécifiques, des abonnements, des services externes ou une
infrastructure de localisation (balises), et le degré de personnalisation dépend de chaque
plateforme.

\subsection{Approches open source et moteurs de routage}

Plusieurs technologies open source permettent de construire des systèmes de routage personnalisés.
Des outils tels que NetworkX, pgRouting, GraphHopper ou OSRM fournissent différents niveaux de
fonctionnalités pour représenter des graphes et calculer des itinéraires. Ces solutions sont
particulièrement intéressantes dans un contexte académique puisqu'elles permettent de conserver un
contrôle important sur la logique de traitement.

\subsection{Comparaison synthétique}

\begin{table}[h]
\centering
\caption{Comparaison synthétique des approches}
\label{tab:compar}
\small
\begin{tabularx}{\linewidth}{|X|X|X|X|}
\hline
\textbf{Approche} & \textbf{Cartographie} & \textbf{Routage} & \textbf{Personnalisation}\\
\hline
Plateformes commerciales & Avancée & Intégré ou disponible & Variable\\
\hline
Moteurs open source & Variable & Avancé & Élevée\\
\hline
Approche personnalisée & Adaptée au besoin & Développé selon le projet & Très élevée\\
\hline
Campus Navigator & Plans SVG dédiés & Graphe personnalisé, guidage temps réel & Adaptée au campus\\
\hline
\end{tabularx}
\end{table}

Cette comparaison montre qu'il existe plusieurs niveaux d'abstraction. Les plateformes commerciales
cherchent à fournir une solution complète, alors qu'une approche personnalisée permet de
construire uniquement les fonctionnalités nécessaires au contexte étudié.

\section{Synthèse et positionnement de Campus Navigator}

L'étude des différentes technologies permet de mieux comprendre les choix effectués pour la
conception de Campus Navigator. Le projet se situe dans une approche de navigation légère et
personnalisée, destinée à un environnement universitaire déterminé, et volontairement sans
infrastructure matérielle.

La représentation des bâtiments sous forme de plans vectoriels fournit une visualisation claire et
adaptée à l'environnement étudié. La structure de graphe représente les relations entre les points
du campus et permet de calculer un itinéraire entre une position de départ et une destination.

L'interface web rend le système accessible à partir d'un simple navigateur, sans installation.
React assure la construction de l'interface et le suivi de la position, tandis que FastAPI fournit
la couche backend.

L'intégration d'un modèle de langage ajoute une dimension conversationnelle : l'utilisateur
formule sa demande de manière naturelle, puis le système identifie la destination correspondante
avant de déclencher le calcul de l'itinéraire. Des mécanismes de repli locaux garantissent que le
service reste disponible même sans le modèle.

Pour la localisation, la solution combine trois sources complémentaires et peu coûteuses : le QR
Code pour le point de départ, le GPS sur les tronçons extérieurs et la podométrie à l'intérieur,
contrainte par l'itinéraire connu.

Cette architecture présente l'avantage de rester relativement simple tout en séparant les
différentes responsabilités du système. Elle constitue une base adaptée à un prototype académique,
tout en laissant la possibilité d'évolutions futures concernant la localisation intérieure, le
stockage des données, la recherche sémantique ou la mise à l'échelle.

\section{Conclusion}

L'état de l'art présenté dans ce chapitre met en évidence le caractère multidisciplinaire de la
navigation intérieure. Une telle solution ne repose pas sur une technologie unique, mais sur
l'intégration de plusieurs composants permettant de représenter l'environnement, déterminer la
position de l'utilisateur, rechercher un itinéraire et présenter le résultat.

Les technologies de localisation telles que le Wi-Fi, le Bluetooth Low Energy et l'UWB offrent
différentes possibilités de positionnement, au prix d'une infrastructure dédiée, tandis que les QR
Codes, le GPS et les capteurs inertiels des smartphones permettent une approche sans matériel
supplémentaire. Sur le plan cartographique, les représentations vectorielles et notamment le format
SVG permettent de construire des plans interactifs adaptés à une application web.

La théorie des graphes fournit un modèle approprié pour représenter les connexions entre les
espaces d'un bâtiment, et les algorithmes tels que Dijkstra et A* permettent d'y rechercher des
chemins optimaux.

L'intelligence artificielle apporte une dimension complémentaire en facilitant l'interaction en
langage naturel. Cette approche doit toutefois être accompagnée de mécanismes de contrôle et de
repli afin de limiter les erreurs d'interprétation et la dépendance à un service externe.

Enfin, l'association de React, Vite, FastAPI, NetworkX, SVG, d'un service d'inférence LLM et d'un
hébergement conteneurisé permet de construire une architecture cohérente pour le prototype
Campus Navigator. Le chapitre suivant présente la conception détaillée du système, son
architecture, ses composants et leur interaction.

% =====================================================================
% Chapitre 3
% =====================================================================
\chapter{Méthodologie et mise en œuvre}

Ce chapitre présente la démarche méthodologique suivie pour la mise en œuvre du système. Il
détaille successivement la mise en place de l'environnement de développement et le choix de la
pile technologique, la conception de la base de données ayant servi à modéliser le campus, le
pipeline de traitement qui relie la requête de l'utilisateur à l'itinéraire, le module de
navigation en temps réel exécuté sur le téléphone, puis le déploiement de l'application.
L'objectif est de présenter non seulement les choix effectués, mais également leur justification
au regard des contraintes du projet.

\section{Mise en place de l'environnement de développement}

Le développement du système repose sur une architecture client-serveur découplée, choisie afin de
séparer clairement la logique métier (identification de la destination, calcul d'itinéraire,
rédaction des instructions) de la couche de présentation (interface conversationnelle, affichage
cartographique, suivi de la position).

Le backend a été développé en Python avec le framework FastAPI, retenu pour sa légèreté, sa
validation native des données (via Pydantic) et sa facilité d'exposition d'endpoints REST.
L'endpoint de navigation est déclaré comme fonction synchrone : FastAPI l'exécute alors dans un
fil d'exécution séparé, ce qui évite que les appels bloquants au modèle de langage ne paralysent
le serveur. Le calcul du plus court chemin repose sur NetworkX, qui permet de modéliser le campus
sous forme de graphe pondéré et d'appliquer directement l'algorithme de Dijkstra.

Le frontend a été développé en React et construit avec Vite. Il est organisé en trois modules : le
composant de chat (\texttt{AdminChat.jsx}), le composant de navigation en temps réel
(\texttt{LiveNavigation.jsx}) et un module de calculs purs, indépendant de React
(\texttt{navMath.js}), regroupant la géométrie le long de l'itinéraire, la détection de pas et les
projections GPS. Pendant le développement, Vite sert l'application en HTTPS (certificat
auto-signé), condition imposée par les téléphones pour accéder au GPS et aux capteurs de
mouvement, et relaie les appels \texttt{/api} vers le backend.

L'intégration de l'intelligence artificielle repose sur l'API Groq, utilisée pour l'extraction de
la destination et pour la rédaction des instructions de navigation. La clé d'accès est fournie par
une variable d'environnement (\texttt{GROQ\_API\_KEY}), jamais versionnée ; en son absence,
l'application fonctionne avec ses mécanismes de repli locaux.

\begin{table}[h]
\centering
\caption{Architecture technique et composants du système}
\label{tab:archi}
\small
\begin{tabularx}{\linewidth}{|p{3.4cm}|p{3.6cm}|X|}
\hline
\textbf{Composant} & \textbf{Technologie} & \textbf{Rôle}\\
\hline
Backend & FastAPI (Python) & Endpoints REST (\texttt{/api/navigate}, \texttt{/api/floorplan}, \texttt{/api/geo}, \texttt{/api/health}), orchestration du pipeline, service du site compilé\\
\hline
Modélisation du graphe & NetworkX & Graphe pondéré du campus, plus court chemin (Dijkstra)\\
\hline
Traitement du langage naturel & Groq (modèle \texttt{openai/gpt-oss-120b}, configurable) & Extraction de la destination, rédaction des instructions de navigation\\
\hline
Repli local & Python (\texttt{unicodedata}, \texttt{difflib}) & Correspondance approximative des lieux, instructions générées par règles\\
\hline
Frontend & React, Vite & Interface conversationnelle, carte interactive, guidage en temps réel\\
\hline
Localisation & API Geolocation, DeviceMotion & GPS sur les tronçons extérieurs, podomètre à l'intérieur\\
\hline
Stockage des données & Excel / pandas, JSON & Données du campus (nœuds, arêtes, lieux, étages), coordonnées GPS des nœuds extérieurs\\
\hline
Génération des QR Codes & Python (\texttt{qrcode}, \texttt{openpyxl}) & Un QR Code par écran, encodant l'adresse publique et le nœud de départ\\
\hline
Édition des plans & Inkscape / SVG & Plans vectoriels des quatre niveaux, support du tracé\\
\hline
Déploiement & Docker, Render & Image unique (site + API) servie en HTTPS\\
\hline
\end{tabularx}
\end{table}

\section{Conception de la base de données}
\label{sec:bdd}

La base de données du campus a été conçue sous forme de classeur Excel (\texttt{Book\_4.xlsx})
structuré en plusieurs feuilles interdépendantes, un choix motivé par la nécessité de disposer
d'un format facilement modifiable et vérifiable manuellement pendant la collecte des données sur le
terrain, avant tout développement logiciel.

Le classeur comprend cinq feuilles :
\begin{itemize}
  \item \textbf{Floors} : 19 niveaux, chacun correspondant au couple (bâtiment, étage). Le campus
  est découpé en 11 bâtiments (bâtiment administratif, atelier, bâtiment de l'amphithéâtre,
  département électrique, amphithéâtres 1 à 4, salle d'étude, service des affaires sociales,
  buvette, sanitaires) et une zone extérieure (\texttt{OUTDOOR}). Chaque ligne précise le numéro
  d'étage et le nombre d'étages du bâtiment (jusqu'à quatre pour le département électrique).
  \item \textbf{Nodes} : 145 points de passage, caractérisés par un identifiant unique
  (\texttt{Node\_Id}), des coordonnées $(x, y)$ relevées sur les plans, un niveau
  d'appartenance (\texttt{floor\_id}) et une description textuelle décrivant un repère visible
  (« en face de la salle B5 », « escalier menant au 1er étage\ldots »). Cette description est
  réutilisée pour rédiger les instructions.
  \item \textbf{Edges} : 174 connexions entre deux nœuds, pondérées par une longueur
  (\texttt{Weight}, interprétée en mètres) et typées (\texttt{Edge\_Type}) : 114 couloirs
  (\texttt{corridor}), 49 allées extérieures (\texttt{outdoor}) et 11 escaliers
  (\texttt{stairs}).
  \item \textbf{Locations} : les 68 lieux nommés du campus, chacun rattaché à un nœud du graphe,
  avec une catégorie et des alias (noms alternatifs séparés par des virgules ou des
  points-virgules). Cette feuille constitue le référentiel interrogé lors de l'extraction de la
  destination.
  \item \textbf{Node\_Types} : la typologie des nœuds : accès à un point d'intérêt
  (\texttt{poi\_access}), intersection, escalier (\texttt{stairs}), point de passage
  (\texttt{waypoint}) et point d'entrée (\texttt{entry\_point}). Les six points d'entrée
  (N3, N77, N87, N88, N114, N142) correspondent aux écrans du campus sur lesquels sont affichés les
  QR Codes.
\end{itemize}

Ce classeur est complété par un fichier \texttt{geo.json} associant à 31 nœuds extérieurs leurs
coordonnées GPS (latitude, longitude), relevées sur le terrain. Ces coordonnées permettent de
suivre l'utilisateur par GPS sur les tronçons extérieurs de l'itinéraire (section~\ref{sec:live}).

\begin{figure}[h]
\centering
\img{0.85\linewidth}{powerbi_modele.png}
\caption{Schéma du modèle de données cartographique extrait de la plateforme Power BI}
\label{fig:modele}
\end{figure}

\begin{figure}[h]
\centering
\img{0.95\linewidth}{excel_nodes.png}
\caption{Extrait de la feuille Nodes, montrant les colonnes Node\_Id, x, y, floor\_id, Description}
\label{fig:nodes}
\end{figure}

Afin de systématiser le contrôle de cette base, un tableau de bord Power BI a été construit,
structuré en deux pages. La première, présentée en Figure~\ref{fig:pbi1}, propose une série
d'indicateurs synthétiques : nombre total de nœuds, d'arêtes et de points d'intérêt, nœuds
orphelins (\emph{Orphan Nodes}), arêtes invalides (\emph{Invalid\_Edges}), coordonnées ou étages
manquants (\emph{Missing\_Coordinates}, \emph{Missing\_Floors}) et doublons d'identifiants
(\emph{Duplicates in Nodes}). La seconde, présentée en Figure~\ref{fig:pbi2}, détaille ces
anomalies ligne par ligne à l'aide de colonnes calculées (\emph{Edge\_Valid},
\emph{From\_Node\_Valid}, \emph{To\_Node\_Valid}, \emph{Self\_Loop},
\emph{Duplicate\_Edge\_Count}), permettant d'identifier précisément quelles arêtes référencent
un nœud inexistant, quelles arêtes forment une boucle, ou quels identifiants sont dupliqués, avant
leur correction dans le fichier source.

\begin{figure}[h]
\centering
\img{0.75\linewidth}{powerbi_qualite.png}
\caption{Vue d'ensemble de la qualité des données (tableau de bord Power BI)}
\label{fig:pbi1}
\end{figure}

\begin{figure}[h]
\centering
\img{0.85\linewidth}{powerbi_integrite.png}
\caption{Contrôle de l'intégrité du graphe (tableau de bord Power BI)}
\label{fig:pbi2}
\end{figure}

Au-delà de ces contrôles structurés, le tableau de bord a également été utilisé de manière
exploratoire, afin de retrouver l'origine d'incohérences repérées lors des tests de l'application.
La Figure~\ref{fig:e67} illustre un tel cas : en filtrant les tables sur les nœuds N88 et
\texttt{To\_Node = N7}, il a été possible d'isoler l'arête à l'origine du problème : E67 (N88
$\rightarrow$ N7, type \texttt{corridor}). Cette arête reliait un nœud intérieur du bâtiment (N88),
sans accès direct vers l'extérieur, à un nœud extérieur (N7) : une connexion qui ne correspond à
aucun passage physique réel. Cette erreur de saisie n'était pas détectable par les contrôles
structurels (les deux nœuds existaient, et aucune boucle ni duplication n'était en cause), mais
l'algorithme de plus court chemin proposait, pour certains trajets, un itinéraire empruntant cette
arête, plaçant l'utilisateur face à une instruction impossible à suivre. Ce cas illustre une limite
propre à ce type de contrôle : la validité structurelle d'un graphe ne garantit pas sa cohérence
sémantique par rapport à la réalité physique qu'il modélise, ce qui a conduit à compléter les
vérifications automatisées par une revue manuelle ciblée des arêtes reliant l'intérieur à
l'extérieur.

\begin{figure}[h]
\centering
\img{0.85\linewidth}{powerbi_e67.png}
\caption{Capture d'écran « Search for Inconsistencies » isolant l'arête E67}
\label{fig:e67}
\end{figure}

Les Figures~\ref{fig:pbi1} à~\ref{fig:e67} correspondent à des captures réalisées à différents
stades du contrôle, avant la correction complète des données. Dans l'état actuel du classeur, le
graphe chargé par NetworkX est connexe : 144 nœuds reliés par 172 connexions distinctes (deux
lignes de la feuille \emph{Edges} décrivant une même connexion sont fusionnées par le graphe non
orienté). Un seul nœud, N129, reste isolé et n'est donc jamais proposé dans un itinéraire.

\section{Pipeline de traitement}

Le fonctionnement du système repose sur un pipeline en six étapes, depuis l'identification du
point de départ jusqu'au guidage de l'utilisateur (Figure~\ref{fig:sequence}).

\begin{enumerate}
  \item \textbf{Identification du point de départ.} L'utilisateur scanne un QR Code affiché sur un
  écran du campus. Ce code encode l'adresse publique de l'application, paramétrée avec
  l'identifiant du nœud correspondant (par exemple \texttt{?start=N114}). Si l'utilisateur
  rescanne un code en cours de trajet, l'application lui propose de reprendre l'itinéraire
  précédent depuis ce nouveau point (la dernière destination est conservée trente minutes dans le
  stockage local du navigateur).

  \item \textbf{Formulation de la requête.} L'utilisateur exprime sa destination en langage
  naturel via l'interface de chat (ex. « Où se trouve la salle B7 ? »), en français ou en anglais.
  La requête part toujours de la position courante de l'utilisateur, mise à jour pendant la marche,
  et non du QR Code initial.

  \item \textbf{Extraction de la destination.} La requête est d'abord transmise au modèle de
  langage, avec la liste des lieux (nom, alias, nœud), en lui demandant de ne répondre que par un
  identifiant de nœud (température nulle) ; l'identifiant est extrait de la réponse par expression
  régulière (\texttt{N\textbackslash d+}) puis vérifié dans le graphe. Si le modèle n'est pas
  disponible ou ne trouve rien, une recherche locale prend le relais :
  \begin{itemize}
    \item normalisation du texte (minuscules, suppression des accents et de la ponctuation) ;
    \item recherche d'un nom ou alias de lieu dont tous les mots significatifs (hors mots vides
    tels que « de », « la », « salle ») figurent dans la requête, en privilégiant la
    correspondance la plus longue ;
    \item à défaut, recherche d'un radical de mot propre à un seul lieu (ex. « bibliothèques »
    $\rightarrow$ « bibliotheque ») ;
    \item enfin, similarité de chaînes (\texttt{difflib}), acceptée uniquement au-delà d'un ratio
    de 0{,}85 pour éviter les correspondances hasardeuses.
  \end{itemize}

  \item \textbf{Calcul de l'itinéraire.} Le nœud de départ et le nœud de destination sont soumis à
  l'algorithme de Dijkstra (\texttt{nx.dijkstra\_path}, pondéré par la longueur des arêtes),
  produisant la séquence ordonnée de nœuds à parcourir. Les cas particuliers renvoient un code
  d'erreur explicite : départ inconnu (\texttt{bad\_start}), destination introuvable
  (\texttt{not\_found}), utilisateur déjà sur place (\texttt{already\_there}) ou absence de chemin
  (\texttt{no\_path}).

  \item \textbf{Génération des instructions.} Le serveur produit \emph{exactement une instruction
  par nœud} du chemin (l'instruction $k$ explique comment atteindre le nœud $k$), ce qui permet au
  frontend d'afficher l'instruction correspondant à la position courante. Pour chaque nœud, il
  construit une ligne de contexte : le lieu (bâtiment et étage), le repère tiré de la description
  du nœud et, le cas échéant, le changement calculé de façon déterministe par rapport au nœud
  précédent (« Sortez du bâtiment », « Entrez dans : \ldots », « Montez l'escalier jusqu'au 1er
  étage »). Ces lignes sont transmises au modèle de langage avec des règles de rédaction
  (impératif, une phrase courte par étape, aucun identifiant technique, ne mentionner un escalier
  à monter que si un changement d'étage est indiqué), en mode de réponse JSON. La réponse est
  ensuite contrôlée : le nombre d'instructions doit être égal au nombre de nœuds, les identifiants
  techniques éventuels sont remplacés, « étage 0 » devient « rez-de-chaussée », et la dernière
  instruction doit annoncer l'arrivée. Si l'un de ces contrôles échoue, ou si le modèle est
  indisponible, des instructions générées par règles, à partir des mêmes lignes de contexte, sont
  utilisées.

  \item \textbf{Restitution et guidage.} Le backend retourne l'itinéraire enrichi de métadonnées
  par nœud ; le frontend l'affiche sur les plans des étages concernés et guide l'utilisateur en
  temps réel (section~\ref{sec:live}).
\end{enumerate}

\begin{figure}[h]
\centering
\resizebox{\linewidth}{!}{%
\begin{tikzpicture}[
  font=\sffamily\small,
  actor/.style={draw, fill=blue!8, rounded corners=2pt, minimum width=2.5cm, minimum height=0.8cm, align=center},
  msg/.style={-{Stealth[length=2.2mm]}, thick},
  ret/.style={-{Stealth[length=2.2mm]}, dashed},
  note/.style={draw, fill=yellow!20, font=\sffamily\scriptsize, align=left, inner sep=2pt},
]
  \tikzset{every node/.append style={fill=white, inner sep=1.5pt}}
  \foreach \name/\x/\label in {u/0/Utilisateur, qr/3/{Écran -- QR Code}, fe/6.5/Frontend (React), be/10.5/Backend (FastAPI), llm/14/API Groq (LLM), nx/17/NetworkX}{
    \node[actor] (\name) at (\x,0) {\label};
    \draw[dashed, gray] (\x,-0.4) -- (\x,-14.2);
  }
  \draw[msg] (0,-1) -- node[above]{1 Scanne le QR Code} (3,-1);
  \draw[msg] (3,-1.7) -- node[above]{2 Ouvre l'URL \texttt{?start=N114}} (6.5,-1.7);
  \draw[msg] (0,-2.6) -- node[above]{3 Formule sa destination} (6.5,-2.6);
  \draw[msg] (6.5,-3.5) -- node[above]{4 POST /api/navigate \{text, startNode, lang\}} (10.5,-3.5);
  \draw[msg] (10.5,-4.4) -- node[above]{5 Extraction de la destination} (14,-4.4);
  \draw[ret] (14,-5.1) -- node[above]{6 Node\_Id (ou repli local)} (10.5,-5.1);
  \draw[msg] (10.5,-6.0) -- node[above, pos=0.6]{7 dijkstra\_path(start, dest)} (17,-6.0);
  \draw[ret] (17,-6.7) -- node[above, pos=0.4]{8 Séquence de nœuds} (10.5,-6.7);
  \draw[msg] (10.5,-7.6) -- node[above]{9 Contexte de chaque nœud} (14,-7.6);
  \draw[ret] (14,-8.3) -- node[above]{10 Instructions (JSON)} (10.5,-8.3);
  \node[note, anchor=west] at (10.7,-9.0) {contrôle du nombre d'étapes,\\repli par règles si besoin};
  \draw[ret] (10.5,-10.0) -- node[above]{11 \{path, pathMeta, edgeMeters, steps\}} (6.5,-10.0);
  \draw[ret] (6.5,-10.9) -- node[above]{12 Trajet sur le plan + instruction} (0,-10.9);
  \draw[msg] (0,-11.8) -- node[above]{13 Marche (GPS / pas détectés)} (6.5,-11.8);
  \draw[ret] (6.5,-12.6) -- node[above]{14 Position, étage et instruction mis à jour} (0,-12.6);
  \draw[msg] (6.5,-13.5) -- node[above]{15 Écart au trajet : POST /api/navigate \{destNode\}} (10.5,-13.5);
\end{tikzpicture}}
\caption{Diagramme de séquence illustrant le pipeline de traitement}
\label{fig:sequence}
\end{figure}

\begin{table}[h]
\centering
\caption{Spécifications des endpoints et paramètres de l'API REST}
\label{tab:api}
\small
\begin{tabularx}{\linewidth}{|p{4.2cm}|p{1.6cm}|X|}
\hline
\textbf{Endpoint} & \textbf{Méthode} & \textbf{Paramètres d'entrée}\\
\hline
\texttt{/api/navigate} & POST & Corps JSON : \texttt{text} (requête en langage naturel), \texttt{startNode} (nœud de départ ou position courante), \texttt{lang} (\texttt{fr} ou \texttt{en}), \texttt{destNode} (optionnel : destination déjà connue, utilisée lors d'un recalcul)\\
\hline
\texttt{/api/floorplan/\{level\}} & GET & Paramètre d'URL : \texttt{level} (\texttt{F0}, \texttt{F1}, \texttt{F2}, \texttt{F3} ; \texttt{OUTDOOR} est ramené à \texttt{F0})\\
\hline
\texttt{/api/geo} & GET & Aucun\\
\hline
\texttt{/api/health} & GET & Aucun\\
\hline
\end{tabularx}
\end{table}

L'endpoint \texttt{/api/navigate} constitue le point d'entrée principal du pipeline. Sa réponse
contient le nom du lieu identifié (\texttt{destinationName}), le nœud de destination
(\texttt{destinationNode}), la séquence ordonnée des nœuds (\texttt{path}), la longueur en mètres
de chaque tronçon (\texttt{edgeMeters}), les instructions (\texttt{steps}, une par nœud) et, pour
chaque nœud, un ensemble de métadonnées (\texttt{pathMeta}) : niveau et plan d'affichage
(\texttt{floorId}, \texttt{mapLevel}), indicateur extérieur, nom du bâtiment, numéro d'étage et
nombre d'étages du bâtiment, coordonnées dans le repère du plan (\texttt{x}, \texttt{y}) et, pour
les nœuds extérieurs, coordonnées GPS (\texttt{geo}). Ces métadonnées permettent au frontend de
dessiner le trajet, de changer de plan et d'afficher les indications d'escalier sans nouvel appel
au serveur. En cas d'échec, la réponse se limite à un champ \texttt{error} accompagné d'un code
(\texttt{code}) traduit par l'interface dans la langue choisie.

Le Listing~\ref{lst:reponse} présente un extrait simplifié de réponse pour la requête « salle B7 »
depuis le nœud N114 : l'itinéraire calculé compte 12 nœuds et 132 mètres, part de l'extérieur,
traverse le bâtiment administratif, le département électrique et le bâtiment de l'amphithéâtre,
puis monte au premier étage de l'atelier.

\begin{lstlisting}[caption={Extrait simplifié d'une réponse de \texttt{/api/navigate}}, label={lst:reponse}]
{
  "destinationName": "Salle B7",
  "destinationNode": "N17",
  "path": ["N114", "N142", "N139", "N138", "N89", "N93",
           "N92", "N91", "N79", "N15", "N16", "N17"],
  "edgeMeters": [8, 7, 40, 5, 5, 23, 18, 10, 5, 6, 5],
  "pathMeta": [
    { "node": "N114", "floorId": "OUTDOOR", "mapLevel": "F0",
      "outdoor": true, "x": 19469.4, "y": 28152.9,
      "geo": [33.9695993, -6.8764596] },
    ...
  ],
  "steps": ["Vous êtes devant ...", "...", "Vous êtes arrivé : Salle B7."]
}
\end{lstlisting}

L'endpoint \texttt{/api/floorplan/\{level\}} répond à un besoin distinct : pour chaque niveau
traversé, le frontend doit afficher le plan correspondant. Plutôt que de transmettre l'ensemble des
plans en une seule fois, cet endpoint fournit à la demande le chemin du fichier SVG (\texttt{file}),
les dimensions de référence du plan (\texttt{viewBox}) et la liste des nœuds du niveau avec leurs
coordonnées converties dans le repère du plan (\texttt{nodes}), cette conversion correspondant à la
calibration décrite au chapitre suivant. Les nœuds extérieurs sont dessinés sur le plan du
rez-de-chaussée, qui représente également les abords des bâtiments. Le frontend met ces réponses
en cache et précharge les plans de tous les niveaux d'un itinéraire, de sorte que le changement
d'étage pendant la marche est instantané.

L'endpoint \texttt{/api/geo} renvoie les nœuds extérieurs géolocalisés (coordonnées GPS et
coordonnées sur le plan) et les allées qui les relient ; il sert à placer un point GPS sur la carte.
Enfin, \texttt{/api/health} indique si le serveur fonctionne, si le modèle de langage est configuré
et combien de nœuds compte le graphe ; il est utilisé par la plateforme d'hébergement pour
surveiller le service.

\section{Module de navigation en temps réel}
\label{sec:live}

Une fois l'itinéraire reçu, le composant de navigation en temps réel prend le relais sur le
téléphone. Il affiche en permanence l'emplacement courant (« Vous êtes ici : bâtiment
administratif · rez-de-chaussée »), l'instruction à suivre, une indication d'escalier lorsque
l'étape suivante change d'étage, une barre de progression, la distance restante et la liste
complète des instructions, dont les étapes franchies sont cochées.

\subsection{Représentation de la position}

La position de l'utilisateur n'est pas représentée par des coordonnées libres, mais par sa
progression le long de l'itinéraire : un couple $(i, p)$ indiquant que l'utilisateur se trouve à la
fraction $p \in [0, 1[$ du tronçon reliant le nœud $i$ au nœud $i+1$. Avancer d'une distance $d$
consiste à consommer la longueur restante des tronçons successifs ; la distance restante jusqu'à
la destination s'écrit :
\begin{equation}
D_{\text{restante}} = \sum_{j=i}^{n-1} \ell_j \; - \; p\,\ell_i
\end{equation}
où $\ell_j$ est la longueur (en mètres) du tronçon $j$ et $n$ le dernier nœud. Sur la carte, le
point est interpolé linéairement entre les nœuds $i$ et $i+1$ lorsqu'ils figurent sur le même
plan. Cette représentation contraint la position à l'itinéraire, ce qui compense en grande partie
l'imprécision des capteurs.

\subsection{Podométrie à l'intérieur}

À l'intérieur, et chaque fois que le GPS n'est pas exploitable, la progression est estimée en
comptant les pas à partir de l'accéléromètre (événement \texttt{devicemotion}). Pour chaque
mesure, on calcule la norme de l'accélération $|a|$, dont on retire la composante de gravité
estimée par un filtre passe-bas, puis on lisse le signal :
\begin{align}
g_t &= 0{,}95\, g_{t-1} + 0{,}05\, |a_t| \\
s_t &= 0{,}6\, s_{t-1} + 0{,}4\, \big(|a_t| - g_t\big)
\end{align}
Un pas est détecté lorsque le signal lissé présente un maximum local supérieur à un seuil
($1{,}2~\text{m/s}^2$), avec une période réfractaire de 300~ms entre deux pas pour éviter les
doubles détections. Chaque pas fait avancer l'utilisateur de 0{,}7~m le long de l'itinéraire.

Sur iOS, l'accès aux capteurs n'est accordé qu'à la suite d'un geste de l'utilisateur : la
demande d'autorisation est donc déclenchée au moment où il appuie sur « Envoyer ». Si aucun
capteur n'est disponible ou si l'accès est refusé, un message l'explique et un mode
« Démo » simulant une marche à 5~m/s reste disponible pour les présentations.

\subsection{Suivi GPS à l'extérieur}

Sur les tronçons dont les deux extrémités possèdent des coordonnées GPS, la position est fournie
par l'API de géolocalisation du navigateur en haute précision. Une mesure n'est exploitée que si
elle date de moins de 8 secondes et que sa précision annoncée est inférieure ou égale à 30 mètres ;
dans le cas contraire, le podomètre reste actif.

Les coordonnées sont converties en mètres dans un repère local centré sur un point de référence
$(\varphi_0, \lambda_0)$, approximation largement suffisante à l'échelle d'un campus :
\begin{equation}
x = R\,(\lambda - \lambda_0)\,\frac{\pi}{180}\,\cos\!\Big(\varphi_0 \frac{\pi}{180}\Big),
\qquad
y = R\,(\varphi - \varphi_0)\,\frac{\pi}{180},
\qquad R = 6\,371\,000~\text{m}
\end{equation}

La mesure est ensuite projetée orthogonalement sur les tronçons extérieurs de l'itinéraire proches
de la position courante (un tronçon en arrière, jusqu'à six en avant). Pour un tronçon $[A, B]$,
avec $\vec{q} = \overrightarrow{AB}$ et $\vec{p} = \overrightarrow{AM}$ le vecteur vers la mesure
$M$ :
\begin{equation}
f = \min\!\Big(\max\!\Big(\frac{\vec{p}\cdot\vec{q}}{\|\vec{q}\|^2}, 0\Big), 1\Big),
\qquad
d_\perp = \big\|\vec{p} - f\,\vec{q}\big\|
\end{equation}
Le tronçon minimisant $d_\perp$ donne la nouvelle progression $(i, f)$. Pour absorber le bruit du
GPS, un recul de moins de 6 mètres le long du trajet est ignoré.

Pour l'affichage, le point GPS n'est pas forcé sur l'itinéraire : il est placé sur l'allée
extérieure la plus proche du réseau complet, en conservant à la fois la position le long de l'allée
et l'écart latéral (converti dans le repère du plan à l'aide du rapport entre la longueur de
l'allée sur le plan et sa longueur en mètres). L'utilisateur voit ainsi ses détours réels, entourés
d'un cercle représentant la précision de la mesure.

\subsection{Détection d'écart et recalcul de l'itinéraire}

L'utilisateur est considéré comme hors itinéraire lorsque sa distance au trajet dépasse
$\max(30~\text{m}, \text{précision})$. Si cette situation dure plus de 6 secondes (ou à la demande,
via un bouton), l'application détermine le nœud extérieur géolocalisé le plus proche. Si ce nœud
fait partie de la suite de l'itinéraire, l'utilisateur a simplement pris un raccourci et la
progression y est avancée ; sinon, un nouvel itinéraire vers la même destination est demandé au
serveur à partir de ce nœud, en transmettant directement le nœud de destination
(\texttt{destNode}) afin d'éviter une nouvelle interprétation de la requête. Un délai minimal de
20 secondes entre deux recalculs évite les recalculs en cascade.

\subsection{Affichage cartographique}

La carte affiche uniquement le plan du niveau où se trouve l'utilisateur et le tronçon de
l'itinéraire situé sur ce plan. Lorsque la progression atteint un nœud situé sur un autre niveau,
le plan change automatiquement ; un marqueur orange indique l'endroit où le trajet quitte le plan
courant, avec une flèche vers le haut ou vers le bas selon le sens de l'escalier, et une épingle
rouge marque la destination lorsqu'elle se trouve sur le plan affiché.

Le trajet restant part de la position de l'utilisateur et raccourcit au fur et à mesure qu'il
avance ; il est tracé en bleu avec un contour blanc et des tirets animés indiquant le sens de
marche, tandis que la portion déjà parcourue est grisée. Par défaut, la caméra suit l'utilisateur
avec un zoom rapproché ; dès qu'il déplace la carte, le suivi est suspendu et un bouton
« Recentrer » apparaît. Les marqueurs et l'épaisseur des traits sont recalculés en fonction du zoom
afin de conserver une taille constante à l'écran, et les déplacements du point sont lissés par une
interpolation exponentielle pour éviter les sauts. Enfin, l'écran est maintenu allumé pendant la
navigation (API \emph{Wake Lock}).

\section{Déploiement}

Dans la version présentée en août, l'application tournait sur l'ordinateur de développement et
les QR Codes pointaient vers son adresse sur le réseau local (\texttt{http://192.168.x.x:5173}),
ce qui ne permettait de l'utiliser qu'aux téléphones connectés au même réseau Wi-Fi. Pour lever
cette contrainte, l'application a été rendue déployable sur une adresse publique.

Le site et l'API sont servis par un seul service : FastAPI sert les fichiers compilés de
l'application React en plus des endpoints \texttt{/api}. Il n'existe ainsi qu'une seule adresse,
sans proxy ni configuration CORS spécifique. L'image Docker est construite en deux étapes : une
image Node.js~20 compile le frontend avec Vite, puis une image Python~3.12 installe les dépendances
du backend et embarque le code du serveur, les données et le site compilé.

L'hébergement utilise la plateforme Render, configurée par un fichier \emph{Blueprint}
(\texttt{render.yaml}) : un service web Docker, surveillé via \texttt{/api/health}, redéployé
automatiquement à chaque envoi sur le dépôt Git, et dont la clé Groq est saisie dans le tableau de
bord de la plateforme. Render fournit automatiquement une adresse HTTPS. Avec l'offre gratuite, le
service se met en veille après environ 15 minutes d'inactivité, et la première requête suivante
peut prendre de 30 à 60 secondes ; une offre payante ou un ping régulier est donc recommandé pour
un usage réel sur le campus.

Enfin, le script de génération des QR Codes lit les points d'entrée dans la feuille
\emph{Node\_Types} et produit, pour chacun, une image encodant \texttt{<adresse
publique>/?start=<Node\_Id>}, avec un niveau de correction d'erreur élevé (Q). Il refuse les
adresses locales (\texttt{localhost}, \texttt{192.168.x.x}) afin d'éviter de reproduire le
problème de réseau, et peut également générer des codes pour d'autres types de nœuds (par exemple
les escaliers), permettant à l'utilisateur de resynchroniser sa position en chemin.

\medskip
Ce chapitre a présenté la méthodologie suivie pour la construction du système, depuis le choix des
outils jusqu'à un pipeline complet reliant la requête de l'utilisateur au calcul de l'itinéraire,
au guidage en temps réel et à la mise en ligne. Une attention particulière a été portée à la
fiabilité des données, contrôlée à l'aide d'un tableau de bord Power BI, ainsi qu'à la robustesse
du service face à l'indisponibilité du modèle de langage. Le chapitre suivant présente les
résultats obtenus ainsi que les difficultés techniques rencontrées et les solutions apportées.

% =====================================================================
% Chapitre 4
% =====================================================================
\chapter{Résultats, difficultés rencontrées et solutions}

Ce chapitre dresse un bilan du travail réalisé. Il présente d'abord les résultats obtenus, sous la
forme de l'application telle qu'elle se présente à l'utilisateur final, puis détaille les
principales difficultés techniques rencontrées au cours du développement — qu'elles soient liées à
des contraintes externes, à la fidélité de la représentation cartographique, à l'ergonomie mobile,
aux capteurs ou au déploiement — ainsi que les solutions mises en œuvre. Cette démarche de
résolution, fondée sur une analyse du code et des données plutôt que sur des corrections
superficielles, constitue une part importante du travail effectué.

\section{Résultats obtenus}

L'application, dans son état actuel, permet à un utilisateur ayant scanné un QR Code de formuler
sa destination en langage naturel et d'obtenir un itinéraire tracé sur les plans réels du campus,
accompagné d'une instruction par étape. Pendant la marche, la carte suit l'utilisateur, le trajet
restant raccourcit, l'instruction courante et la distance restante se mettent à jour, et le plan
change automatiquement lorsque le trajet passe à un autre étage. L'interface est disponible en
français et en anglais, gère les trajets traversant plusieurs bâtiments et plusieurs étages, et
reste utilisable aussi bien sur ordinateur que sur smartphone. Déployée sur une adresse HTTPS
publique, elle est accessible depuis n'importe quel réseau (Wi-Fi ou données mobiles).

Au niveau mobile :

\begin{figure}[h]
\centering
\begin{subfigure}[t]{0.47\linewidth}
  \centering
  \img{\linewidth}{mobile_carte.png}
  \caption{Carte et instruction courante}
  \label{fig:mob1}
\end{subfigure}\hfill
\begin{subfigure}[t]{0.47\linewidth}
  \centering
  \img{\linewidth}{mobile_instructions.png}
  \caption{Liste des instructions générées}
  \label{fig:mob2}
\end{subfigure}
\caption{Interface mobile lors de la recherche de la Salle B7}
\label{fig:mobile}
\end{figure}

Transition d'un itinéraire multi-niveaux du rez-de-chaussée vers le premier étage :

\begin{figure}[h]
\centering
\begin{subfigure}[t]{0.47\linewidth}
  \centering
  \img{\linewidth}{mobile_escalier.png}
  \caption{Indication d'escalier et marqueur de sortie du plan au rez-de-chaussée}
  \label{fig:mob3}
\end{subfigure}\hfill
\begin{subfigure}[t]{0.47\linewidth}
  \centering
  \img{\linewidth}{mobile_etage1.png}
  \caption{Changement automatique vers le plan du 1er étage}
  \label{fig:mob4}
\end{subfigure}
\caption{Changement d'étage automatique au cours d'un itinéraire}
\label{fig:etages}
\end{figure}

\begin{figure}[h]
\centering
\begin{subfigure}[t]{0.47\linewidth}
  \centering
  \img{\linewidth}{mobile_gps.png}
  \caption{Suivi GPS sur un tronçon extérieur, avec cercle de précision}
  \label{fig:mob5}
\end{subfigure}\hfill
\begin{subfigure}[t]{0.47\linewidth}
  \centering
  \img{\linewidth}{mobile_recalcul.png}
  \caption{Itinéraire recalculé après un écart au trajet}
  \label{fig:mob6}
\end{subfigure}
\caption{Guidage en extérieur : suivi GPS et recalcul automatique}
\label{fig:gps}
\end{figure}

La Figure~\ref{fig:gps} illustre le guidage sur la partie extérieure du campus, par exemple vers la
buvette des enseignants, accessible uniquement par une porte extérieure. Sur ces tronçons, le point
bleu suit la position GPS réelle ; lorsque l'utilisateur s'éloigne durablement du trajet,
l'itinéraire est recalculé depuis le point le plus proche et un message l'en informe.

\section{Difficultés rencontrées et solutions}

Plusieurs difficultés techniques significatives ont jalonné le développement du projet. Elles sont
présentées ci-après par ordre chronologique approximatif d'apparition.

\subsection{Contraintes liées au modèle de langage}
\label{sec:llm}

L'implémentation initiale du module d'intelligence artificielle reposait sur l'API Gemini
(bibliothèque \texttt{google-genai}). Elle s'est heurtée à des limites de quota bloquantes dès les
premiers tests : le service retournait une erreur 429, avec une limite fixée à 0 requête pour le
modèle utilisé, rendant tout appel impossible dans le cadre d'un usage gratuit. Une migration vers
l'API Groq a été effectuée, celle-ci offrant un niveau gratuit compatible avec les besoins du
développement ; le modèle \emph{llama-3.1-8b-instant} a été retenu dans un premier temps.

Quelques jours après l'envoi du rapport d'avancement du 3 août 2026, ce modèle a cessé de répondre
de manière fiable. L'application a alors été migrée vers \emph{openai/gpt-oss-20b}, puis vers
\emph{openai/gpt-oss-120b}, plus fiable dans le respect des consignes de rédaction, sans
modification de l'architecture d'appel : le modèle est désormais un simple paramètre de
configuration (\texttt{GROQ\_MODEL}).

Cet épisode illustre une contrainte propre à l'usage de modèles tiers : leur disponibilité et leurs
conditions d'accès évoluent indépendamment du code applicatif. Il a conduit à deux choix
structurants. D'une part, une intégration faiblement couplée au modèle exact utilisé. D'autre part,
la mise en place de mécanismes de repli pour chacune des deux tâches confiées au modèle :
correspondance approximative locale pour l'identification de la destination, et instructions
générées par règles pour le guidage. Le service reste ainsi utilisable même sans clé d'API ou en
cas de panne du fournisseur.

Une seconde difficulté a concerné la qualité des instructions générées. Les premières versions
recevaient la suite brute des descriptions de nœuds et produisaient parfois des instructions
incohérentes : nombre d'étapes différent du nombre de points de passage (rendant impossible
l'affichage de l'instruction correspondant à la position), consigne de monter un escalier
simplement cité comme repère, identifiants techniques (« N109 ») recopiés dans le texte, ou
« étage 0 » au lieu de « rez-de-chaussée ». La solution a consisté à calculer de façon
déterministe, côté serveur, les informations critiques (bâtiment, étage, entrée, sortie, sens de
l'escalier) et à les transmettre explicitement au modèle, à lui imposer une réponse au format JSON,
puis à valider systématiquement cette réponse avant de l'utiliser.

\emph{N.B. : l'erreur 429 signifie que trop de requêtes ont été envoyées au serveur dans un laps de
temps trop court.}

\subsection{Constitution et fiabilité de la base de données}

La constitution de la base de données du campus a été l'une des tâches les plus longues du projet.
Le recensement des 68 points d'intérêt et la saisie des 145 nœuds et 174 arêtes ont nécessité un
travail manuel important : relevé des coordonnées, retranscription dans le classeur Excel, puis
vérification croisée avec les plans en cours de conception. Le suivi GPS a ajouté une étape
supplémentaire : le relevé sur le terrain des coordonnées GPS des 31 nœuds extérieurs.

Cette phase s'est révélée intrinsèquement itérative : chaque test de l'application pouvait révéler
une incohérence — nœud mal rattaché à son étage, arête manquante ou mal typée, coordonnée erronée —
nécessitant un retour au fichier source, une correction, puis un nouveau test. Le tableau de bord
Power BI présenté en section~\ref{sec:bdd} a été conçu pour systématiser cette vérification. Ce
travail de fiabilisation reste néanmoins continu : de nouvelles incohérences ont été identifiées
plusieurs semaines après la mise en place initiale du graphe, et un nœud (N129) reste aujourd'hui
isolé. Cela souligne la nécessité de maintenir ces contrôles tout au long du développement plutôt
que de les considérer comme une étape ponctuelle.

\subsection{Choix de l'outil de conception des plans}

La conception des plans vectoriels a fait l'objet d'une réorientation en cours de projet. Une
première version a été réalisée avec draw.io, choisi pour sa simplicité de prise en main. Cette
approche s'est révélée insuffisante pour représenter fidèlement l'agencement réel des bâtiments :
le niveau de détail atteignable restait limité, et la structure des fichiers SVG exportés se
prêtait mal à un traitement précis des coordonnées. Les plans ont donc été entièrement repris sous
Inkscape, un éditeur vectoriel offrant un contrôle plus fin sur le tracé et les coordonnées. Cette
réorientation a impliqué une phase d'apprentissage de l'outil, jusqu'alors non maîtrisé par
l'équipe, avant de pouvoir reproduire l'ensemble des plans avec la précision requise — précision
indispensable à la calibration présentée ci-après.

\subsection{Calibration cartographique des plans}

Une fois les plans redessinés sous Inkscape, une difficulté est apparue au moment de superposer
l'itinéraire calculé : la ligne représentant le trajet n'apparaissait pas à l'endroit attendu,
bien que le calcul du plus court chemin soit correct côté serveur.

L'origine du problème résidait dans une confusion entre deux systèmes de coordonnées présents dans
un même fichier SVG : le \texttt{viewBox}, repère interne utilisé par le navigateur pour positionner
les éléments du plan, et le repère de la page affiché par la barre d'outils d'Inkscape (exprimé en
points, pt), depuis lequel les coordonnées des nœuds avaient été relevées puis reportées dans le
classeur Excel. Ces deux repères, bien que définis dans le même fichier, ne coïncident pas.

La calibration a été résolue en établissant le facteur de conversion exact entre ces deux repères,
déduit des dimensions déclarées par les fichiers SVG (\texttt{viewBox} de largeur 39\,841,941 pour
une page de 4\,870~pt), puis validée par comparaison avec des nœuds de référence positionnés
manuellement :
\begin{equation}
\text{échelle} = \frac{39\,841{,}941}{4\,870} = 8{,}181097
\end{equation}
\begin{equation}
X_{\text{réel}} = 8{,}181097 \times X_{\text{pt}} - 4000, \qquad
Y_{\text{réel}} = 8{,}181097 \times Y_{\text{pt}} - 4000
\end{equation}

La première relation détermine le rapport d'échelle entre le repère vectoriel du fichier et les
mesures relevées sous Inkscape, garantissant l'homothétie de la transformation. Les deux suivantes
appliquent ce facteur aux coordonnées initiales, en introduisant une translation constante de
$-4000$ unités qui corrige le décalage d'origine du \texttt{viewBox} ($-4000$, $-4000$). Cette
conversion, appliquée côté serveur avant tout envoi au frontend, est commune aux quatre plans.

\begin{figure}[h]
\centering
\img{0.95\linewidth}{calibration_avant_apres.png}
\caption{Comparaison avant/après de l'alignement des nœuds sur un même étage (version desktop)}
\label{fig:calib}
\end{figure}

Cette phase a également révélé une anomalie : un itinéraire traversant plusieurs bâtiments situés
au même étage affichait un onglet distinct pour chacun, alors qu'ils correspondent au même plan
physique. La correction a consisté à regrouper les tronçons de l'itinéraire par plan d'affichage
plutôt que par bâtiment : chaque nœud reçoit un niveau de plan (\texttt{F0} à \texttt{F3}, les
nœuds extérieurs étant rattachés au plan du rez-de-chaussée). Ce même principe a ensuite permis de
remplacer les onglets par un changement de plan automatique, piloté par la position de
l'utilisateur.

\subsection{Ergonomie mobile}

Le comportement de la carte interactive sur smartphone a posé plusieurs difficultés. Dans une
première version, le zoom et le déplacement, fonctionnels sur ordinateur, étaient inopérants sur
mobile : leur gestion reposait exclusivement sur des événements de souris (\texttt{mousedown},
\texttt{mousemove}, \texttt{wheel}), absents sur les terminaux tactiles. Par ailleurs, la
désactivation du comportement tactile par défaut du navigateur — nécessaire pour un contrôle
personnalisé du zoom — avait pour effet, en l'absence de gestionnaire tactile, de bloquer toute
interaction.

Une gestion tactile dédiée a d'abord été ajoutée (déplacement à un doigt, pincement à deux doigts).
Lors de la refonte de la carte pour le guidage en temps réel, la gestion des gestes a été unifiée à
l'aide des \emph{Pointer Events}, qui représentent de la même manière la souris, le doigt et le
stylet : un pointeur actif déclenche un déplacement, deux pointeurs déclenchent un zoom
proportionnel à l'écartement des doigts, et la capture du pointeur garantit que le geste se
poursuit même lorsque le doigt sort de la carte. Le gestionnaire de la molette est déclaré
explicitement comme non passif, certains navigateurs ignorant silencieusement l'annulation du
défilement de la page dans le cas contraire.

L'interface a également été repensée pour les petits écrans : mise en page occupant toute la
hauteur de l'écran (y compris les zones réservées des téléphones à encoche), champs de saisie
suffisamment grands pour éviter le zoom automatique d'iOS, boutons adaptés au toucher et logo de
l'ENSAM dans l'en-tête. La lisibilité du trajet a été améliorée (tracé élargi, contour contrasté,
marqueurs distincts pour la position, la destination et les changements d'étage).

\subsection{Accès depuis le réseau du campus}

Lors des tests de la version d'août, l'application n'était accessible qu'aux téléphones connectés
au même réseau Wi-Fi que l'ordinateur de développement : les QR Codes encodaient l'adresse locale
de cet ordinateur (\texttt{http://192.168.x.x:5173}), et le backend y tournait également. Un
utilisateur en données mobiles, ou sur un autre réseau, ne pouvait pas ouvrir l'application.

La solution a consisté à déployer l'application sur une adresse publique (section précédente) et à
régénérer les QR Codes à partir de cette adresse. Le script de génération refuse désormais les
adresses locales, afin que ce problème ne puisse pas être réintroduit par erreur. Le regroupement
du site et de l'API derrière une seule adresse a en outre supprimé toute question de proxy ou de
partage de ressources entre origines.

Lors de la mise en ligne, il est également apparu que la clé de l'API Groq avait été enregistrée
dans l'historique du dépôt Git. Elle a été révoquée et remplacée, et la clé est désormais fournie
exclusivement par une variable d'environnement, exclue du dépôt.

\subsection{Accès aux capteurs et précision de la localisation}

Le suivi de la position a soulevé trois difficultés. Premièrement, les navigateurs mobiles
n'autorisent l'accès au GPS et aux capteurs de mouvement que sur des pages sécurisées : servie en
HTTP sur le réseau local, l'application ne recevait aucune mesure. Le serveur de développement a
donc été configuré en HTTPS (certificat auto-signé), et l'hébergement fournit un certificat valide.

Deuxièmement, iOS n'accorde l'accès à l'accéléromètre qu'à la suite d'un geste de l'utilisateur, et
uniquement si la demande est faite immédiatement, avant toute opération asynchrone. La demande
d'autorisation a donc été placée au début du traitement de l'appui sur « Envoyer », le module de
navigation attendant ensuite la réponse avant d'écouter les capteurs.

Troisièmement, aucune des deux sources n'est fiable seule : le podomètre dérive (longueur de pas
fixe, pas parfois manqués), et le GPS est bruité, imprécis à proximité des bâtiments et inutilisable
à l'intérieur. La solution retenue combine les deux sources selon le contexte (GPS lorsque la mesure
est récente et précise sur un tronçon extérieur, podomètre sinon), contraint la position à
l'itinéraire, ignore les petits reculs dus au bruit du GPS et n'engage un recalcul qu'après un écart
prolongé. Lorsque l'utilisateur rescanne un QR Code en chemin, sa position est recalée exactement.

\medskip
Ce chapitre a présenté les résultats obtenus ainsi que les principales difficultés rencontrées lors
de leur mise en œuvre. Au-delà de leur diversité apparente — fiabilité des données, calibration
cartographique, ergonomie mobile, accès réseau, capteurs —, ces difficultés partagent un point
commun : elles n'ont, pour la plupart, été révélées ni par la conception initiale ni par les
contrôles automatisés, mais par les tests successifs de l'application dans des conditions proches
de l'usage réel. L'arête reliant à tort un nœud intérieur à un nœud extérieur, l'écart de
calibration entre le repère d'Inkscape et celui du navigateur, l'absence de gestion tactile, les QR
Codes limités au réseau local ou les capteurs bloqués hors HTTPS sont autant d'anomalies qui
n'affectaient en rien la validité structurelle du système, mais compromettaient son usage concret.
Cette phase de test et de correction itérative, bien que chronophage, a constitué une part
essentielle du travail réalisé.

% =====================================================================
% Conclusion générale
% =====================================================================
\chapter*{Conclusion générale}
\addcontentsline{toc}{chapter}{Conclusion générale}

Ce projet avait pour objectif de concevoir un système de navigation permettant aux étudiants, aux
enseignants et aux visiteurs de l'ENSAM de Rabat de trouver facilement leur chemin au sein du
campus en formulant leur destination en langage naturel. Cet objectif a été atteint sous la forme
d'un prototype fonctionnel et accessible en ligne : après avoir scanné un QR Code identifiant son
point de départ, l'utilisateur formule sa demande dans une interface de type chat, obtient un
itinéraire tracé sur les plans réels du campus, puis est guidé pas à pas, en français ou en anglais,
la carte suivant sa progression et changeant d'étage automatiquement.

Sur le plan technique, le travail a d'abord consisté à modéliser le campus : les 68 points
d'intérêt, les 145 nœuds et les 174 arêtes du graphe ont été collectés et structurés dans un
classeur Excel, puis contrôlés à l'aide d'un tableau de bord Power BI. Le graphe est exploité avec
NetworkX et le plus court chemin calculé par l'algorithme de Dijkstra. Un backend FastAPI orchestre
l'ensemble du pipeline : identification de la destination par un modèle de langage via l'API Groq,
avec une correspondance approximative locale en repli, calcul de l'itinéraire, puis rédaction d'une
instruction par point de passage, contrôlée avant utilisation. Ce choix d'architecture confie au
modèle de langage uniquement l'interprétation et la rédaction du texte, tandis que le calcul du
chemin et les informations critiques (étage, escalier, entrée, sortie) restent déterminés par un
traitement déterministe et vérifiable. Côté client, une interface React affiche le trajet sur les
plans SVG conçus sous Inkscape et suit l'utilisateur grâce au GPS à l'extérieur et à un podomètre
fondé sur l'accéléromètre à l'intérieur, avec recalcul automatique de l'itinéraire en cas d'écart.
L'ensemble est conteneurisé avec Docker et déployé en HTTPS.

Le développement a été l'occasion de rencontrer et de résoudre des difficultés concrètes :
l'indisponibilité de modèles de langage liée aux quotas, qui a conduit à une intégration faiblement
couplée et à des mécanismes de repli ; la calibration entre le repère d'Inkscape et le
\texttt{viewBox} du navigateur, résolue par un facteur de conversion exact ; la gestion des gestes
sur mobile ; l'accès limité au réseau local, résolu par le déploiement public ; les restrictions
d'accès aux capteurs ; ou encore une arête erronée reliant un nœud intérieur à un nœud extérieur,
qui a montré qu'un graphe structurellement valide n'est pas nécessairement cohérent avec la réalité
physique qu'il représente. Ces anomalies ont pour la plupart été révélées par des tests en
conditions réelles, ce qui souligne l'importance d'une démarche itérative de test et de correction.

Le système présente néanmoins des limites propres à un prototype. À l'intérieur, la position est
estimée par podométrie, avec une longueur de pas fixe : la précision se dégrade sur les longs
trajets et dépend de la manière dont le téléphone est tenu. Le GPS n'est exploité que sur les
31 nœuds extérieurs géolocalisés. Les données sont stockées dans un fichier Excel, adapté à cette
phase mais peu approprié à une modification collaborative. L'hébergement gratuit impose un temps de
réveil après une période d'inactivité. Enfin, la qualité des instructions dépend en partie d'un
service externe, même si l'application reste fonctionnelle sans lui.

Plusieurs perspectives d'amélioration peuvent être envisagées. Sur le plan de la localisation, des
QR Codes placés aux escaliers et aux carrefours (déjà pris en charge par le script de génération)
permettraient de recaler régulièrement la position ; l'utilisation de balises Bluetooth Low Energy,
via une application native, ou la calibration de la longueur de pas par utilisateur amélioreraient
la précision à l'intérieur. Sur le plan des données, une migration vers PostgreSQL/PostGIS et une
interface d'administration faciliteraient la mise à jour des lieux et des plans. La recherche de
destinations pourrait être enrichie par des techniques sémantiques (embeddings). Le calcul
d'itinéraires pourrait intégrer des critères d'accessibilité (ascenseurs, rampes) et être étendu à
l'ensemble des bâtiments de l'établissement, et l'interface pourrait être complétée par d'autres
langues, dont l'arabe. Enfin, des tests à plus grande échelle auprès des étudiants permettraient de
valider l'ergonomie auprès d'un public plus large.

Sur le plan personnel et académique, ce projet a permis de mobiliser et de relier des compétences
issues de la filière INDIA : traitement et fiabilisation des données, théorie des graphes,
développement web, exploitation de capteurs, intégration de modèles de langage, conception
d'interfaces et déploiement. Il a surtout offert l'occasion de conduire un projet de bout en bout,
de la collecte des données sur le terrain jusqu'à une application en ligne, en apprenant à
diagnostiquer et à corriger méthodiquement des problèmes concrets.

% =====================================================================
% Glossaire
% =====================================================================
\chapter*{Glossaire}
\addcontentsline{toc}{chapter}{Glossaire}

\begin{description}[leftmargin=0pt, itemsep=0.3em]
\item[A*] : Algorithme de recherche de chemin qui combine le coût déjà parcouru et une estimation du
coût restant (heuristique) pour explorer le graphe de façon plus ciblée que Dijkstra.
\item[Accéléromètre] : Capteur du téléphone mesurant les accélérations selon trois axes ; utilisé
ici pour détecter les pas de l'utilisateur.
\item[API (Application Programming Interface)] : Interface permettant à deux applications de
communiquer selon des règles définies ; utilisée pour les échanges entre le frontend et le backend,
avec le service de modèle de langage et avec les capteurs du navigateur.
\item[Arête] : Dans un graphe, connexion reliant deux nœuds ; représente ici un passage physique
entre deux points du campus (couloir, escalier, allée extérieure).
\item[Backend] : Partie du système exécutée côté serveur, responsable de la logique métier.
\item[BLE (Bluetooth Low Energy)] : Variante basse consommation du Bluetooth, utilisée avec des
balises (\emph{beacons}) pour la localisation intérieure.
\item[Boucle (self-loop)] : Arête reliant un nœud à lui-même, considérée comme une anomalie.
\item[Calibration] : Processus consistant à faire correspondre les coordonnées des nœuds du graphe
avec leur position réelle sur les plans affichés.
\item[Dijkstra] : Algorithme de calcul du plus court chemin entre deux nœuds d'un graphe pondéré,
utilisé ici pour déterminer l'itinéraire optimal.
\item[Docker] : Outil de conteneurisation permettant d'empaqueter une application et ses
dépendances dans une image exécutable sur n'importe quel serveur.
\item[Embeddings] : Représentations numériques de textes permettant de comparer leur sens, utilisées
dans la recherche sémantique.
\item[Endpoint] : Point d'accès d'une API REST, identifié par une URL.
\item[Erreur 429] : Code HTTP indiquant que trop de requêtes ont été envoyées à un serveur dans un
temps donné (limite de quota dépassée).
\item[FastAPI] : Framework Python utilisé pour développer le backend de l'application.
\item[Fingerprinting (empreinte Wi-Fi)] : Technique de localisation consistant à comparer les
signaux Wi-Fi captés avec une cartographie préalablement enregistrée.
\item[Frontend] : Partie du système exécutée côté client, correspondant à l'interface visible.
\item[GNSS / GPS] : Systèmes de positionnement par satellites, efficaces en extérieur mais peu
fiables à l'intérieur des bâtiments.
\item[Graphe] : Structure de données composée de nœuds et d'arêtes, utilisée pour modéliser le
réseau de passages du campus.
\item[Groq] : Service d'inférence donnant accès, via une API, à des modèles de langage hébergés
(ici successivement llama-3.1-8b-instant, openai/gpt-oss-20b puis openai/gpt-oss-120b).
\item[HTTPS] : Version sécurisée du protocole HTTP, exigée par les navigateurs pour accéder au GPS et
aux capteurs de mouvement.
\item[Heuristique] : Estimation du coût restant jusqu'à la destination, utilisée par l'algorithme A*.
\item[Homothétie] : Transformation géométrique qui multiplie toutes les distances par un même
facteur ; utilisée ici pour convertir les coordonnées entre deux repères.
\item[Inkscape] : Éditeur de graphisme vectoriel utilisé pour concevoir les plans des étages au
format SVG.
\item[JSON (JavaScript Object Notation)] : Format de données structuré, utilisé pour les échanges
entre le frontend et le backend et pour les réponses du modèle de langage.
\item[LLM (Large Language Model)] : Modèle de langage de grande taille, utilisé pour l'extraction de
la destination et la rédaction des instructions.
\item[Mécanisme de repli (fallback)] : Traitement de secours utilisé lorsque le traitement principal
(ici le modèle de langage) est indisponible ou renvoie une réponse invalide.
\item[Nœud] : Dans un graphe, point représentant un emplacement précis du campus (intersection,
accès à une salle, escalier, point d'entrée).
\item[Nœud orphelin] : Nœud qui n'est relié à aucune arête, détecté lors du contrôle de qualité.
\item[Node.js] : Environnement d'exécution JavaScript utilisé pour compiler le frontend.
\item[NetworkX] : Bibliothèque Python de manipulation de graphes, utilisée pour la modélisation du
campus et le calcul du plus court chemin.
\item[pandas] : Bibliothèque Python de manipulation de données tabulaires, utilisée pour charger le
classeur Excel.
\item[Panoramique (pan)] : Déplacement du plan à l'écran, à la souris ou avec un doigt.
\item[PDR (Pedestrian Dead Reckoning)] : Navigation pédestre à l'estime, consistant à estimer le
déplacement à partir des pas détectés par les capteurs inertiels.
\item[Pinch-to-zoom] : Geste tactile à deux doigts permettant de zoomer ou dézoomer.
\item[Pipeline] : Enchaînement structuré des étapes de traitement d'une requête.
\item[Pointer Events] : Événements du navigateur représentant de manière unifiée la souris, le
toucher et le stylet.
\item[Power BI] : Outil d'analyse et de visualisation de données de Microsoft, utilisé pour le
contrôle de la cohérence de la base de données.
\item[Projection orthogonale] : Opération consistant à trouver le point d'un segment le plus proche
d'un point donné ; utilisée pour placer une mesure GPS sur l'itinéraire.
\item[Pydantic] : Bibliothèque Python de validation des données, intégrée à FastAPI.
\item[QR Code] : Code-barres bidimensionnel encodant une information, utilisé ici pour identifier le
point de départ de l'utilisateur.
\item[React] : Bibliothèque JavaScript utilisée pour le développement du frontend.
\item[Recalcul d'itinéraire] : Nouveau calcul du plus court chemin déclenché lorsque l'utilisateur
s'écarte durablement du trajet.
\item[Recherche sémantique] : Recherche fondée sur le sens d'une requête plutôt que sur la
correspondance exacte des mots.
\item[Render] : Plateforme d'hébergement utilisée pour déployer l'application.
\item[REST (Representational State Transfer)] : Style d'architecture pour la conception d'API,
reposant sur des requêtes HTTP standard.
\item[SIG (Système d'Information Géographique)] : Système de stockage, d'analyse et de
représentation de données géographiques.
\item[SVG (Scalable Vector Graphics)] : Format d'image vectorielle utilisé pour les plans du campus.
\item[TALN (Traitement Automatique du Langage Naturel)] : Domaine de l'intelligence artificielle
consacré à l'interprétation et à la génération de texte en langage humain.
\item[UWB (Ultra-Wideband)] : Technologie radio permettant de mesurer des distances avec une grande
précision, au prix d'un matériel spécialisé.
\item[ViewBox] : Attribut d'un fichier SVG définissant le repère de coordonnées interne utilisé pour
positionner les éléments du dessin.
\item[Vite] : Outil de développement et de construction d'applications web, utilisé avec React.
\item[PostgreSQL / PostGIS] : Système de gestion de base de données et son extension pour les
données géographiques, envisagés comme évolution du stockage.
\item[Wake Lock] : API du navigateur permettant d'empêcher la mise en veille de l'écran.
\item[UX (User Experience)] : Expérience utilisateur, c'est-à-dire la qualité globale de
l'interaction avec l'application.
\end{description}

% =====================================================================
% Bibliographie
% =====================================================================
\begin{thebibliography}{99}
\addcontentsline{toc}{chapter}{Bibliographie}

\bibitem{modo} Modo Team, « Campus navigation is more than maps: How indoor and outdoor wayfinding
improves the student experience », \emph{Modo Labs}, 16 juill. 2026. [En ligne]. Disponible :
\url{https://modolabs.com/industry/higher-education/campus-navigation/}.

\bibitem{lakshmi} K. Muthu Lakshmi, R. Manojkumar, R. Kesavan et S. Arul Kumaran, « Campus route
map », \emph{Int. Res. J. Adv. Eng. Hub (IRJAEH)}, vol.~3, n\textsuperscript{o}~3, p.~496--501,
mars 2025, e-ISSN 2584-2137, doi : 10.47392/IRJAEH.2025.0069.

\bibitem{hafashimana} E. Hafashimana, « Conception et mise en œuvre d'un système de navigation
intérieure économique utilisant des technologies open-source dans les bâtiments universitaires »,
Mémoire de M.Sc., Univ. du Québec à Chicoutimi, Chicoutimi, QC, Canada, 2024. [En ligne].
Disponible : \url{https://constellation.uqac.ca/id/eprint/10118/1/Hafashimana_uqac_0862N_11282.pdf}.

\bibitem{kunhoth} J. Kunhoth, A. G. Karkar, S. Al-Maadeed et A. Al-Ali, « Indoor positioning and
wayfinding systems: a survey », \emph{Human-centric Comput. Inf. Sci.}, vol.~10,
art.~n\textsuperscript{o}~18, 2020, doi : 10.1186/s13673-020-00222-0.

\bibitem{dedes} G. Dedes et A. G. Dempster, « Indoor GPS positioning: challenges and
opportunities », dans \emph{Proc. IEEE 62nd Veh. Technol. Conf. (VTC-2005-Fall)}, Dallas, TX, USA,
sept. 2005, doi : 10.1109/VETECF.2005.1557943.

\bibitem{polestar} « Bluetooth Low Energy positioning. Room \& zone accuracy, at scale. »,
\emph{Polestar}. [En ligne]. Disponible : \url{https://www.polestar.eu/technologies/ble-positioning/}.

\bibitem{kim} D. Kim, J.-H. Park et Y.-J. Suh, « A Wi-Fi fingerprinting indoor localization
framework using feature-level augmentation via variational graph auto-encoder »,
\emph{Electronics}, vol.~14, n\textsuperscript{o}~14, art.~n\textsuperscript{o}~2807, juill. 2025,
doi : 10.3390/electronics14142807.

\bibitem{gnss} « The difference between GNSS and GPS explained », \emph{Global GPS Systems},
3 nov. 2020 (mis à jour en mai 2026). [En ligne]. Disponible :
\url{https://globalgpssystems.com/gnss/the-difference-between-gnss-and-gps-explained/}.

\bibitem{uwb} « Ultra-wideband (UWB) », dans \emph{RTLS Glossary}, Locaxion. [En ligne]. Disponible :
\url{https://locaxion.com/resources/rtls-glossary/ultra-wideband-uwb/}.

\bibitem{qr} Shreesh, « En quoi les QR Codes figurant sur les plans du campus facilitent-ils
l'orientation ? », \emph{The QR Code Generator} (blog), mis à jour le 4 déc. 2025. [En ligne].
Disponible : \url{https://www.the-qrcode-generator.com/fr/blog/campus-map-qr-code}.

\bibitem{leedev} LeeDev428, « From Static Floor Plans to Searchable Campus Navigation with SVG and
A* », \emph{LeeDev} (blog), 2026. [En ligne]. Disponible :
\url{https://leedev.vercel.app/blog/svg-campus-navigation-a-star-pathfinding}.

\bibitem{inkscape} K. Schnider, « Inkscape : éditeur SVG, dessin vectoriel », \emph{Libellules}
(blog), 31 janv. 2024. [En ligne]. Disponible :
\url{https://www.libellules.net/post/inkscape-\%C3\%A9diteur-svg-dessin-vectoriel}.

\bibitem{dijkstra} « Dijkstra's algorithm », \emph{Coddy}. [En ligne]. Disponible :
\url{https://coddy.tech/visualize/fr/graphs/dijkstras-algorithm}.

\bibitem{astar} R. Kumar, « L'algorithme A* : un guide complet », \emph{DataCamp}, févr. 2025.
[En ligne]. Disponible : \url{https://www.datacamp.com/fr/tutorial/a-star-algorithm}.

\bibitem{nx1} NetworkX Developers, « NetworkX: Network analysis in Python ». [En ligne].
Disponible : \url{https://networkx.org/en/}.

\bibitem{nx2} NetworkX Developers, « Tutorial », dans \emph{NetworkX Documentation}, version stable.
[En ligne]. Disponible : \url{https://networkx.org/documentation/stable/tutorial.html}.

\bibitem{fastapi} S. Ramírez, « FastAPI documentation ». [En ligne]. Disponible :
\url{https://fastapi.tiangolo.com/}.

\bibitem{groq} Groq, « Supported Models », \emph{GroqCloud Documentation}. [En ligne]. Disponible :
\url{https://console.groq.com/docs/models}.

\bibitem{mdngeo} MDN Web Docs, « Geolocation API », Mozilla. [En ligne]. Disponible :
\url{https://developer.mozilla.org/en-US/docs/Web/API/Geolocation_API}.

\bibitem{mdnmotion} MDN Web Docs, « DeviceMotionEvent », Mozilla. [En ligne]. Disponible :
\url{https://developer.mozilla.org/en-US/docs/Web/API/DeviceMotionEvent}.

\bibitem{mdnpointer} MDN Web Docs, « Pointer events », Mozilla. [En ligne]. Disponible :
\url{https://developer.mozilla.org/en-US/docs/Web/API/Pointer_events}.

\bibitem{mdnwake} MDN Web Docs, « Screen Wake Lock API », Mozilla. [En ligne]. Disponible :
\url{https://developer.mozilla.org/en-US/docs/Web/API/Screen_Wake_Lock_API}.

\bibitem{difflib} Python Software Foundation, « difflib — Helpers for computing deltas »,
\emph{Python 3 Documentation}. [En ligne]. Disponible :
\url{https://docs.python.org/3/library/difflib.html}.

\bibitem{docker} Docker Inc., « Multi-stage builds », \emph{Docker Docs}. [En ligne]. Disponible :
\url{https://docs.docker.com/build/building/multi-stage/}.

\bibitem{render} Render, « Blueprint YAML Reference », \emph{Render Docs}. [En ligne]. Disponible :
\url{https://render.com/docs/blueprint-spec}.

\end{thebibliography}

\end{document}
